\chapter{First Probability}\label{ch:probability}

\begin{center}\itshape
$n$ guests leave their hats at the door and get them back at random. The probability that nobody gets their own hat\\
tends to $\frac1e\approx0.368$, whether there are ten guests or ten million.
\end{center}

Ten guests, or ten million. The chance that every hat is wrong settles near $1/e$. The number of guests has almost dropped out. That is a different kind of sentence from the sum of cubes in Chapter~\ref{ch:sums}. There both sides were equal for every $n$, by an identity. Here a sequence of finite probabilities approaches a limit. By the end of the chapter you will be able to say which sentence you have proved.

The subject began with a stake. In the summer of 1654 the Chevalier de M\'er\'e asked Blaise Pascal how to divide the pot of a game that had stopped early. Suppose one player still needs two points and the other needs three, and each point is a fair coin. Fermat's count finishes the story in imagination. Four further games are enough, because $2+3-1=4$. If those four games are equally likely to go either way, there are $2^4=16$ possible records. The player who needs two points fails only in the records where he wins fewer than two of them: $\binom40+\binom41=5$ records. The other $11$ are his. A fair division of the pot is $11/16$ against $5/16$. Some of the four games would not have been played, because the match would already have ended. The winner is the same on every way of filling out a record, so the extra games do not change the share. Pascal reached the same numbers by dividing the stake one game at a time. Their letters run from July to October 1654.

\portraits{pascal}{Blaise Pascal (1623--1662). Anonymous copy of a portrait by Fran\c cois Quesnel, c.\,1690, Palace of Versailles. Public domain, via Wikimedia Commons.}{fermat}{Pierre de Fermat (c.\,1607--1665; the birth year is disputed). Portrait, seventeenth century, artist unknown. Public domain, via Wikimedia Commons.}

Christiaan Huygens saw the correspondence in Paris and wrote the first printed book of the subject. It appeared in 1657 as an appendix to Frans van Schooten's \emph{Exercitationes mathematicae}. The tract itself occupies pages 521 to 534: fourteen propositions, and five problems left to the reader. Several of them divide a stake. Gerolamo Cardano had written on dice in the sixteenth century; that book was printed only in 1663. Galileo had already counted three dice. A total of $9$ can happen in $25$ ways, a total of $10$ in $27$. The totals are not the outcomes. The outcomes are the triples, and there are $6^3=216$ of them.

In 1708 Pierre R\'emond de Montmort asked a question that sounds like a card game. Shuffle cards numbered $1$ to $n$ and count how many land in their own place. What is the chance that none does? The 1713 edition of his book prints his correspondence with Nicolaus Bernoulli, Jacob Bernoulli's nephew. The proportion barely moves once $n$ is past $4$: for five cards it is $44/120=0.3666\ldots$, already within $0.002$ of $1/e$. Euler took up the same \emph{jeu de rencontre} in a memoir printed in 1753 and wrote the alternating sum you will meet below. The counting chapter already has the exact formula. This chapter turns the formula into a probability, and then into a limit.

The same year 1713 saw Jacob Bernoulli's \emph{Ars Conjectandi}, published eight years after his death. The theorem he called golden says that the frequency of an event, in many independent trials, settles near its probability. We prove that statement at the end of the chapter, as the weak law of large numbers. Chebyshev's inequality of 1867 makes the proof four lines long. Ir\'en\'ee-Jules Bienaym\'e had published the variance bound in 1853. The sources record both papers. The argument below follows Chebyshev's use of it.

\portraits{montmort}{Pierre R\'emond de Montmort (1678--1719). Portrait by Alexis-Simon Belle, about 1715. Public domain, via Wikimedia Commons.}{jbernoulli}{Jacob Bernoulli (1655--1705). Portrait by his brother Nicolaus Bernoulli, 1687. Public domain, via Wikimedia Commons.}

The surprises are concrete, and each one is a count with a model written first. Twenty-three people already make a shared birthday more likely than not. Richard von Mises published that threshold in 1939. Figure~\ref{fig:birthday} draws the curve. The bound is the estimate proved below, and Exercise~\ref{ex:birthday} asks you to write that comparison. Collecting all $n$ coupons takes, on average, $nH_n$ boxes, about $n\ln n$. That expectation needs a waiting time, so it waits for Section~\ref{sec:countable-probability} and Exercise~\ref{ex:coupon}. Any $n^2+1$ distinct real numbers contain a monotone subsequence of length $n+1$. That count is Exercise~\ref{ex:erdos} in Chapter~\ref{ch:counting}. In the three-door game it pays to switch, once you have written what the host knows. Marilyn vos Savant said so in her column of 9 September 1990. She later estimated about $10{,}000$ letters, many from people with doctorates, telling her she was wrong. Steve Selvin had posed the game in 1975. The letters do not decide it. Bayes' formula does, and Exercise~\ref{ex:monty} shows that a different host gives a different answer.

Finite experiments come first. A short countable extension is marked as such. One probability space that holds an infinite sequence of trials waits for Volume~4.

\section{Finite experiments}

A probability model starts with a nonempty set $\Omega$ of possible outcomes. You choose that set before you compute. For one die, $\Omega=\{1,\ldots,6\}$. For two labelled dice it is $\{1,\ldots,6\}^2$, thirty-six ordered pairs. An \emph{event} is a subset of $\Omega$. ``The sum is seven'' is the set of six pairs in Figure~\ref{fig:dice}. The possible sums $2,\ldots,12$ are another description of the same rolls. They are not equally likely, because they do not contain the same number of pairs.

\begin{figure}[ht]
\centering
\fig{ch09-dice}
\caption{Each cell is one ordered outcome of two fair dice, so each has mass $1/36$. The six rust cells are the outcomes that sum to $7$. The charcoal cell is the only outcome that sums to $2$. Seven is six times as likely as two because six pairs add to seven and one pair is two ones.}
\label{fig:dice}
\end{figure}

For finite $\Omega$, assign each outcome a \emph{mass} $p(\omega)\ge0$ with $\sum_{\omega\in\Omega}p(\omega)=1$, and put
\[
 P(A)=\sum_{\omega\in A}p(\omega)\qquad(A\subseteq\Omega).
\]
Probability is a weighted count. It satisfies $P(\Omega)=1$ and $P(A\cup B)=P(A)+P(B)$ when $A\cap B=\varnothing$. These two rules, together with nonnegativity, determine the masses the other way: $p(\omega)=P(\{\omega\})$, and finite additivity rebuilds the sum. Also $P(\varnothing)=0$, since adding the empty event changes nothing. These are the \emph{finite probability axioms}.

The \emph{uniform} model gives each outcome mass $1/|\Omega|$, hence $P(A)=|A|/|\Omega|$. Uniformity is an assumption about the experiment. A fair die justifies it. A shuffled pack justifies it, if the shuffle really is uniform. A list of names for the outcomes does not. On two fair dice, each ordered pair has mass $1/36$, so the sum is seven with probability $6/36$ and the sum is two with probability $1/36$.

\begin{think}
Two fair coins are tossed. Is the list ``zero heads, one head, two heads'' a uniform sample space? Write the four elementary outcomes before you answer.
\end{think}

The four outcomes are HH, HT, TH, TT. If the coins are fair and independent, each has mass $1/4$. ``One head'' contains two of them. The three-word list is a legal set of events. It is not a uniform set of outcomes.

The same count, with the inequality $1-x\le e^{-x}$, answers the birthday question. Take $365$ days and $n$ people, with birthdays uniform and independent. The probability that all $n$ birthdays are different is
\[
 \prod_{k=1}^{n-1}\Bigl(1-\frac{k}{365}\Bigr),
\]
for $n\le365$, and it is $0$ if $n>365$. Proposition~\ref{prop:convex} gives $1-x\le e^{-x}$ for real $x$, because $e^{-x}$ is convex and its tangent at $0$ is $1-x$. The product is therefore at most
\[
 \exp\Bigl(-\sum_{k=1}^{n-1}\frac{k}{365}\Bigr)=\exp\Bigl(-\frac{n(n-1)}{730}\Bigr).
\]
For $n=23$, $23\cdot22/730=253/365$, and the bound is $\exp(-253/365)$. The exponent is $-253/365$. Compare $253/365$ with $\ln2$. Chapter~\ref{ch:sequences} gives $\ln2<0.693148$. The product $693148\cdot365=252999020$ is less than $253\cdot10^6$, so $0.693148<253/365$ and $\ln2<253/365$. The exponent $-253/365$ is then less than $\ln(1/2)$, and $\exp(-253/365)<1/2$. The probability that all twenty-three birthdays differ is less than $1/2$. A shared birthday is then more likely than not. The exact probability that all twenty-three birthdays differ is about $0.4927$, so the probability of a shared birthday is about $0.5073$. Exercise~\ref{ex:birthday} asks you to write the comparison. Figure~\ref{fig:birthday} shows why $23$ feels early: the curve is the pairs, and there are $\binom{23}2=253$ of them.

\begin{figure}[ht]
\centering
\fig{ch08-birthday}
\caption{Left: the probability that among $n$ people two share a birthday (charcoal, exact), with the bound $1-e^{-n(n-1)/730}$ of Exercise~\ref{ex:birthday} (dashed). It crosses $1/2$ at $n=23$ and exceeds $0.99$ at $n=57$. Right: the expected number of boxes in the coupon collector's problem of Exercise~\ref{ex:coupon}, $nH_n$ (rust), against the simulated average of $2000$ runs (dots); the dashed curve is $n\ln n$.}
\label{fig:birthday}
\end{figure}

\section{Conditioning and independence}

If we learn that an event $B$ has occurred, and $P(B)>0$, the outcomes inside $B$ keep their relative weights. We rescale so that those weights add to one. That is the definition
\[
 P(A\mid B)=\frac{P(A\cap B)}{P(B)}.
\]
It is a probability on the restricted sample space $B$. Nonnegativity, total mass one, and additivity on disjoint events follow from the same properties of $P$. The definition does not assign $P(A\mid B)$ when $P(B)=0$.

Check it on a die. Let $B$ be the even faces, so $B=\{2,4,6\}$ and $P(B)=1/2$. Let $A$ be the event that the face exceeds $3$. Then $A\cap B=\{4,6\}$, and
\[
 P(A\mid B)=\frac{2/6}{3/6}=\frac23.
\]
Two of the three surviving faces lie above $3$. Figure~\ref{fig:even} is that sentence as a picture.

\begin{figure}[ht]
\centering
\fig{ch09-even}
\caption{A fair die, after we are told that the face is even. The faces $1$, $3$ and $5$ are out. The three that remain share the conditional mass. Two of them, marked $>3$, exceed $3$, so the conditional probability is $2/3$.}
\label{fig:even}
\end{figure}

Events $A$ and $B$ are \emph{independent} when $P(A\cap B)=P(A)P(B)$. If $P(B)>0$, this says exactly that $P(A\mid B)=P(A)$: learning $B$ did not move the probability of $A$. Disjointness is a different assertion. Two disjoint events of positive probability cannot be independent, because their intersection has probability zero, and the product of two positive numbers is positive. An event of probability zero is independent of every event, by the product rule, even though conditioning on it is undefined.

A family $A_1,\ldots,A_r$ is \emph{mutually independent} when
\[
 P\Bigl(\bigcap_{i\in I}A_i\Bigr)=\prod_{i\in I}P(A_i)
\]
for every nonempty subset $I\subseteq\{1,\ldots,r\}$. Asking this only for pairs is \emph{pairwise independence}. Figure~\ref{fig:pairsigns}, in the section on variance, is a case where every pair is independent and the triple is not.

For independent repetition of a finite experiment with masses $p(e)$ on a set $E$, use $\Omega=E^n$ and
\[
 P(\{(e_1,\ldots,e_n)\})=\prod_{i=1}^n p(e_i).
\]
The masses add to $\bigl(\sum_{e\in E}p(e)\bigr)^n=1$. If the event in coordinate $i$ restricts that coordinate to $A_i\subseteq E$, distributing the finite sums gives
\[
 P(e_i\in A_i\text{ for every }i)=\prod_i\sum_{e\in A_i}p(e).
\]
The trials are independent because the model says so. Repeating an experiment does not, by itself, make the copies independent. You have to put the product masses in.

Two rules follow at once and get used constantly. First, $P(A^c)=1-P(A)$. Second, the \emph{law of total probability}: if $B_1,\ldots,B_k$ partition $\Omega$ into events of positive probability, then
\[
 P(A)=\sum_i P(A\mid B_i)\,P(B_i).
\]
Split $A$ along the $B_i$ and apply the definition of conditional probability. Bayes' formula is the revision written as a quotient. If $P(A)>0$ and $P(B)>0$,
\[
 P(B\mid A)=\frac{P(A\mid B)\,P(B)}{P(A)}.
\]
Both sides equal $P(A\cap B)/P(A)$. A hypothesis $B$ had probability $P(B)$. You observe $A$. The new probability of the hypothesis is the left-hand side. Proposition~\ref{prop:proba} records the formula again, next to the rules for averages. The game show is the test. Write the host's rule before you compute.

\begin{examplex}[Monty Hall]
Three doors hide one car and two goats. You pick door $1$. The host, who knows where the car is, opens one of the other two doors to reveal a goat, choosing at random if both hide goats, and offers you the remaining closed door. Should you switch? Let $C_i$ be ``the car is behind door $i$'', each of probability $1/3$, and let $H_3$ be ``the host opens door $3$''. Then $P(H_3\mid C_1)=1/2$, because he chooses between doors $2$ and $3$. Next, $P(H_3\mid C_2)=1$, because door $2$ hides the car and he must open door $3$. And $P(H_3\mid C_3)=0$. Total probability gives
\[
 P(H_3)=\frac13\cdot\frac12+\frac13\cdot1+0=\frac12.
\]
Bayes' formula gives
\[
P(C_2\mid H_3)=\frac{P(H_3\mid C_2)P(C_2)}{P(H_3)}=\frac{1\cdot\frac13}{\frac12}=\frac23,
\qquad
P(C_1\mid H_3)=\frac{\frac12\cdot\frac13}{\frac12}=\frac13.
\]
Switching wins with probability $2/3$. Figure~\ref{fig:monty} is the same arithmetic as a tree. The guess ``two doors left, so fifty-fifty'' forgets that the host was not allowed to open the door with the car. That constraint carries information. Exercise~\ref{ex:monty} changes the host. The answer moves with him, which is why the argument has to be conditional probability and not a feeling about the two remaining doors.
\end{examplex}

\begin{figure}[ht]
\centering
\fig{ch08-monty}
\caption{The Monty Hall game as a tree, when you first pick door $1$. Each branch carries its probability. The leaves give the probability of each pair (car, opened door). Given that door $3$ is opened, the shaded leaves have total probability $1/2$, and the car is behind door $2$ with conditional probability $(1/3)/(1/2)=2/3$. The same tree with ``the host opens door $2$'' gives the same conclusion by symmetry. Switching wins two times in three.}
\label{fig:monty}
\end{figure}

A test with a high accuracy can still raise mostly false alarms, if the condition is rare. The denominator of Bayes' formula is where the count of those alarms sits.

\begin{examplex}[a test in a population of ten thousand]
Suppose a uniformly chosen person has a condition with probability $0.01$. A test is positive with probability $0.99$ when the condition is present, and with probability $0.05$ when it is absent. Let $D$ mean ``condition'' and $+$ mean ``positive''. Then
\[
P(+)=0.99\cdot0.01+0.05\cdot0.99=0.0594,
\qquad P(D\mid+)=\frac{0.0099}{0.0594}=\frac16.
\]
In a population of ten thousand, the same rates give $99$ true positives and $495$ false positives. Figure~\ref{fig:bayes} draws those two counts. They are $10000$ times the probabilities, not a guaranteed total in one town. Neither the sensitivity $0.99$ nor the specificity $0.95$ is the probability of having the condition after a positive result. That probability is $1/6$.
\end{examplex}

\begin{figure}[ht]\centering\fig{review08-bayes}
\caption{Positive results per ten thousand people, counted at the rates of the example. The pale bar is $99$ true positives. The dark bar is $495$ false positives. The two fills differ in gray, so the comparison does not depend on colour.}\label{fig:bayes}\end{figure}

\subsection*{Does a second test settle it?}
\textbf{Mission.} Suppose a second test has the same rates, and the two tests are independent \emph{conditional on} whether the condition is present. Find $P(D\mid+,+)$. What changes if the second test simply repeats the first result?

\textbf{Hint.} Square the conditional rates in the first model. Do not square the posterior $1/6$.

\textbf{Solution.} Under conditional independence,
\[
P(D\mid+,+)=\frac{0.01(0.99)^2}{0.01(0.99)^2+0.99(0.05)^2}
=\frac{99}{124}\approx0.7984.
\]
If the second result is a copy of the first, then two positives are the same event as one positive, and the answer remains $1/6$. The same marginal accuracies can yield very different posteriors. Independence has to describe the mechanism. It is not a factor you insert because multiplication looks convenient.

\section{Expectation and the hat check}

A \emph{random variable} is a function $X:\Omega\to\R$. The notation $\{X=x\}$ means the inverse image $X^{-1}(\{x\})$. Likewise $\{X\ge a\}$ is an event. The \emph{law}, or distribution, of $X$ lists the values and the masses $P(X=x)$. Those events partition $\Omega$, so the masses add to one. The \emph{expectation} is the sum over outcomes,
\[
 \mathbb E[X]=\sum_{\omega\in\Omega}X(\omega)\,p(\omega).
\]
It is a number computed from the model, not a forecast of one outcome. The \emph{variance} and the \emph{standard deviation} are
\[
 \operatorname{Var}X=\mathbb E[(X-\mathbb EX)^2],\qquad
 \sigma(X)=\sqrt{\operatorname{Var}X}.
\]
Variance measures squared distance from the mean.

\begin{proposition}[first rules]\label{prop:proba}
\begin{enumerate}[label=(\alph*)]
\item $P(A)=\sum_{\omega\in A}P(\{\omega\})$, and \textbf{Bayes}: $P(B\mid A)=\frac{P(A\mid B)P(B)}{P(A)}$ when $P(A),P(B)>0$.
\item $\mathbb E[X]=\sum_xx\,P(X=x)$, the sum over the values of $X$.
\item \textbf{Linearity}: $\mathbb E[X+Y]=\mathbb E[X]+\mathbb E[Y]$ and $\mathbb E[cX]=c\,\mathbb E[X]$, with \emph{no independence needed}. For an event $A$, $\mathbb E[\mathbf 1_A]=P(A)$.
\item $\operatorname{Var}X=\mathbb E[X^2]-\mathbb E[X]^2$.
\end{enumerate}
\end{proposition}
\begin{proof}
(a) Additivity, by induction on $|A|$. Bayes: both sides equal $P(A\cap B)/P(A)$.
(b) Group the outcomes $\omega$ by the value $x=X(\omega)$, and use (a).
(c) The definition of $\mathbb E$ is a sum over $\omega$ of $X(\omega)P(\{\omega\})$, linear in $X$. And $\sum_\omega\mathbf 1_A(\omega)P(\{\omega\})=P(A)$.
(d) Expand $(X-m)^2$ with $m=\mathbb EX$ and use (c): $\mathbb E[X^2]-2m^2+m^2$.
\end{proof}

Part (b) is how you compute from the law. A fair die has values $1,\ldots,6$, each of mass $1/6$, so its expectation is $(1+\cdots+6)/6=7/2$. The die has never shown $3.5$, and it is not about to. The average belongs to the six masses, not to a face of the cube.

For the number $X$ of heads in two fair tosses, the values $0,1,2$ have probabilities $1/4$, $1/2$, $1/4$. Part (b) gives $\mathbb EX=1$. You can see the same number without listing the law. Let $I_1$ and $I_2$ be the indicators of heads on the two tosses. Then $X=I_1+I_2$, and part (c) gives $\mathbb EX=1/2+1/2=1$. The indicators are functions on the same $\Omega$. Linearity does not ask them to be independent.

Variance is nonnegative, and it is zero exactly when $X=\mathbb EX$ at every outcome of positive mass: a sum of nonnegative terms is zero exactly when every positive-weight term is zero. For the two-coin count,
\[
 \operatorname{Var}X=\frac14\cdot1^2+\frac12\cdot0^2+\frac14\cdot1^2=\frac12.
\]
Part (d) gives the same number as $\mathbb E[X^2]-1$. The indicator $\mathbf1_A$ takes the value $1$ on $A$ and $0$ elsewhere. Part (c) says $\mathbb E[\mathbf1_A]=P(A)$. A count of events is a sum of indicators. Its expectation is the sum of the probabilities, whether or not the events overlap.

\begin{think}
On average, how many guests get their own hat back when there are $10$ guests? And with a million guests? Guess first, then read the example.
\end{think}

\begin{example}
In a uniformly random permutation of $\{1,\ldots,n\}$, with $n\ge1$, the expected number of fixed points is
\[
 \sum_{i=1}^n P(\sigma(i)=i)=n\cdot\frac1n=1.
\]
\end{example}

The law of the number of fixed points is awkward. Figure~\ref{fig:derang}, on the right, draws it for $n=10$. Nobody would compute the expectation from that law. Linearity never builds the law. The indicators $\mathbf1_{\sigma(i)=i}$ are not independent, and the expectation does not care. The average number of correctly returned hats is exactly $1$, for ten guests and for a million. The party does not get more fixed points, on average, by getting larger. It gets more ways to be tangled, and those ways cancel in the expectation. Exercises~\ref{ex:fixedvar}, \ref{ex:cycles}, \ref{ex:coupon} and~\ref{ex:records} are the same trick on four different counts.

The probability of no fixed point is a different sentence. Chapter~\ref{ch:counting} proved
\[
 \frac{D_n}{n!}=\sum_{k=0}^n\frac{(-1)^k}{k!}.
\]
Taylor's formula for the exponential, Exercise~\ref{ex:taylorln}, identifies the infinite sum with $e^{-1}$. The terms $1/k!$ decrease to $0$, so the alternating-series bound of Proposition~\ref{prop:signedseries} gives
\[
 \Bigl|\frac{D_n}{n!}-\frac1e\Bigr|\le\frac1{(n+1)!}.
\]
For $n=10$ the error is at most $1/11!$, which is about $2.51\times10^{-8}$, and the true gap is about $2.31\times10^{-8}$. Figure~\ref{fig:derang} shows the oscillation. Multiply the bound by $n!$. For $n\ge2$,
\[
 \Bigl|D_n-\frac{n!}{e}\Bigr|\le\frac1{n+1}<\frac12,
\]
so $D_n$ is the integer nearest to $n!/e$. The case $n=1$ is separate: $D_1=0$, and the distance from $1!/e$ to $0$ is $1/e$. The series for $e$ gives $e>1+1+1/2=5/2>2$, hence $1/e<1/2$. The integer $1$ is farther from $1/e$ than $0$ is. Thus $D_n$ is the integer nearest to $n!/e$ for every $n\ge1$. The first eight values are
\[
 D_1,\ldots,D_8=0,1,2,9,44,265,1854,14833.
\]
The recurrence $D_n=(n-1)(D_{n-1}+D_{n-2})$ is Exercise~\ref{ex:derangrec}. The counting chapter records the two forms Euler wrote.

\begin{figure}[ht]
\centering
\fig{ch08-derangements}
\caption{Left: the proportion $D_n/n!$ of permutations with no fixed point, for $n=1,\ldots,10$. It oscillates about $1/e$ and sits about $2.3\times10^{-8}$ away at $n=10$. Right: the law of the number of fixed points of a random permutation of $10$ elements (charcoal), computed as $\binom{10}{k}D_{10-k}/10!$, against the Poisson masses $e^{-1}/k!$ (coral). The bars agree to four decimals. The proof that the laws converge is in Volume~4.}
\label{fig:derang}
\end{figure}

\section{Variance and the weak law}

Expectation says where a random quantity sits on average. Variance says how far from there it typically lands. The theorems in this section make that precise: first for a sum, then as an inequality that uses nothing but the mean and the variance, then as Bernoulli's golden theorem. None of them needs more than finite sums.

\begin{definitionx}[independent random variables, covariance]
Random variables $X$ and $Y$ on $\Omega$ are \emph{independent} if $P(X=x,\,Y=y)=P(X=x)P(Y=y)$ for all values $x,y$. A family $X_1,\ldots,X_n$ is \emph{pairwise independent} if every pair is independent, and \emph{independent} if $P(X_1=x_1,\ldots,X_n=x_n)=\prod_i P(X_i=x_i)$ for all values $x_i$. The \emph{covariance} is
\[
 \operatorname{Cov}(X,Y)=\mathbb E[(X-\mathbb EX)(Y-\mathbb EY)]=\mathbb E[XY]-\mathbb E[X]\,\mathbb E[Y],
\]
the second form by expanding, as for the variance. In particular $\operatorname{Var}X=\operatorname{Cov}(X,X)$.
\end{definitionx}

Pairwise independence does not give independence of the whole family. Take two independent fair signs $A,B\in\{-1,1\}$, and set $X=A$, $Y=B$, $Z=AB$. Each sign is fair. Each pair $(X,Y)$, $(X,Z)$, $(Y,Z)$ takes its four possible values with probability $1/4$ each. But $XYZ=1$ at every outcome. In particular $P(X=Y=Z=1)=1/4$, whereas $P(X=1)P(Y=1)P(Z=1)=1/8$. Figure~\ref{fig:pairsigns} is the four-row table. The covariance identities below need only pairwise independence. A product law for three or more coordinates needs the joint hypothesis.

\begin{figure}[ht]
\centering
\fig{ch09-signs}
\caption{Two independent fair signs $A$ and $B$, and $Z=AB$. Each row has mass $1/4$. The pairs of columns $(A,B)$, $(A,Z)$ and $(B,Z)$ are each uniform on their four sign pairs, so $X=A$, $Y=B$ and $Z$ are pairwise independent. The fourth column records the mass of each row. The rust row has mass $1/4$. The product of the three separate probabilities of $+1$ is $1/8$. Pairwise independence is not independence of the triple.}
\label{fig:pairsigns}
\end{figure}

\begin{proposition}[products and sums]\label{prop:varsum}
\begin{enumerate}[label=(\alph*)]
\item If $X$ and $Y$ are independent, $\mathbb E[XY]=\mathbb E[X]\,\mathbb E[Y]$, that is, $\operatorname{Cov}(X,Y)=0$. The converse is false.
\item $\operatorname{Var}(aX+b)=a^2\operatorname{Var}X$, and $\operatorname{Var}(X+Y)=\operatorname{Var}X+\operatorname{Var}Y+2\operatorname{Cov}(X,Y)$.
\item If $X_1,\ldots,X_n$ are pairwise independent, $\operatorname{Var}(X_1+\cdots+X_n)=\operatorname{Var}X_1+\cdots+\operatorname{Var}X_n$.
\end{enumerate}
\end{proposition}
\begin{proof}
(a) Group the outcomes by the pair of values, as in Proposition~\ref{prop:proba}(b):
\[
\begin{aligned}
\mathbb E[XY]
&=\sum_{x,y}xy\,P(X=x,Y=y)
=\sum_{x,y}xy\,P(X=x)P(Y=y)\\
&=\Bigl(\sum_x x\,P(X=x)\Bigr)\Bigl(\sum_y y\,P(Y=y)\Bigr).
\end{aligned}
\]
For the converse, let $X$ be uniform on $\{-1,0,1\}$ and let $Y=X^2$. Then $\mathbb E[XY]=\mathbb E[X^3]=0=\mathbb E[X]\,\mathbb E[Y]$, but
\[
 P(X=0,Y=0)=\frac13\ne P(X=0)P(Y=0)=\frac13\cdot\frac13.
\]

(b) The centred variable $aX+b-\mathbb E[aX+b]$ equals $a(X-\mathbb EX)$, and squaring gives the first claim. For the second, write $X'=X-\mathbb EX$ and $Y'=Y-\mathbb EY$. Then
\[
 \operatorname{Var}(X+Y)=\mathbb E[(X'+Y')^2]=\mathbb E[X'^2]+\mathbb E[Y'^2]+2\mathbb E[X'Y'].
\]

(c) Induct on $n$, using (b). The step is
\[
 \operatorname{Var}(S_{n-1}+X_n)=\operatorname{Var}S_{n-1}+\operatorname{Var}X_n+2\operatorname{Cov}(S_{n-1},X_n),
\]
and $\operatorname{Cov}(S_{n-1},X_n)=\sum_{i<n}\operatorname{Cov}(X_i,X_n)=0$ by (a). Covariance is linear in each argument, because it is a difference of expectations, and expectations are linear.
\end{proof}

Covariance records how two centred quantities vary together. For random variables $X_1,\ldots,X_d$, the \emph{covariance matrix} $\Sigma$ is the symmetric matrix with entries $\Sigma_{ij}=\operatorname{Cov}(X_i,X_j)$. Take any column of real coefficients $u=(u_1,\ldots,u_d)^{\mathsf T}$. Expanding the square gives
\[
 u^{\mathsf T}\Sigma u
 =\mathbb E\Bigl[\Bigl(\sum_i u_i(X_i-\mathbb EX_i)\Bigr)^2\Bigr]
 =\operatorname{Var}\Bigl(\sum_i u_i X_i\Bigr)\ge0.
\]
Every covariance matrix is positive semidefinite, in the sense of Chapter~\ref{ch:linear}. The diagonal entries are variances. An off-diagonal entry can be negative. If $X$ is a fair zero-or-one variable and $Y=1-X$, then
\[
 \Sigma=\frac14\begin{pmatrix}1&-1\\-1&1\end{pmatrix}.
\]
The eigenvalues are $0$ and $1/2$. The direction $(1,1)$ has variance zero because $X+Y=1$ is fixed. The perpendicular direction measures the random difference. A zero eigenvalue need not be a defect in the data. It can be an exact relation.

The spectral theorem supplies an orthogonal matrix $Q$ such that $Q^{\mathsf T}\Sigma Q$ is diagonal. The centred vector $Z=Q^{\mathsf T}(X-\mathbb EX)$, with the coordinate means collected in $\mathbb EX$, has this diagonal covariance, because each entry follows by distributing finite sums. Its coordinates are uncorrelated. They need not be independent. The pair $X$ and $X^2$ in the proof of Proposition~\ref{prop:varsum} is already a warning. Exercise~\ref{ex:covariance-sampling} finds the directions for a sample of fixed size, without expanding a determinant.

Variance adds for pairwise independent variables. The sum of $n$ independent copies of $X$ has mean $n\,\mathbb EX$ and standard deviation $\sqrt n\,\sigma(X)$. The ratio of spread to mean therefore shrinks like $1/\sqrt n$. That scale is the binomial clause of the next proposition. The factor $(N-n)/(N-1)$ is the correction when the draws are dependent, because they are taken without replacement.

\begin{proposition}[the classical laws]\label{prop:laws}
\begin{enumerate}[label=(\alph*)]
\item (\textbf{Bernoulli}) Let $0\le p\le1$. If $P(X=1)=p$ and $P(X=0)=1-p$, then $\mathbb E X=p$ and $\operatorname{Var}X=p(1-p)\le1/4$.
\item (\textbf{binomial}) Let $n\ge0$ be an integer and $0\le p\le1$, and let $S$ be the number of successes in $n$ independent trials, each succeeding with probability $p$. Then
\[
 P(S=k)=\binom nk p^k(1-p)^{n-k}\qquad(0\le k\le n),
\]
and $\mathbb ES=np$, $\operatorname{Var}S=np(1-p)$.
\item (\textbf{hypergeometric}) Let $N\ge2$, $0\le K\le N$ and $0\le n\le N$ be integers. Draw $n$ objects without replacement from $N$ objects of which $K$ are marked, and let $X$ be the number of marked objects drawn. Then
\[
 P(X=k)=\frac{\binom Kk\binom{N-K}{n-k}}{\binom Nn},
 \qquad
 \mathbb E X=n\frac KN,
\]
and
\[
 \operatorname{Var}X=n\frac KN\Bigl(1-\frac KN\Bigr)\frac{N-n}{N-1}.
\]
\item (\textbf{uniform}) For an integer $n\ge1$, if $X$ is uniform on $\{1,\ldots,n\}$, then $\mathbb EX=(n+1)/2$ and $\operatorname{Var}X=(n^2-1)/12$.
\end{enumerate}
\end{proposition}
\begin{proof}
(a) $\mathbb EX=p$ and $\mathbb E[X^2]=p$, so $\operatorname{Var}X=p-p^2$. The quadratic $p(1-p)$ has maximum $1/4$ at $p=1/2$.

(b) For $p=0$ or $p=1$ the count is deterministic. The displayed formula uses the empty-product convention $0^0=1$. For $0<p<1$, work in the product model $\Omega=\{0,1\}^n$. A particular sequence with $k$ successes has probability $p^k(1-p)^{n-k}$, and there are $\binom nk$ such sequences. Now $S=X_1+\cdots+X_n$, where $X_i$ is the indicator of success at trial $i$. The $X_i$ are independent, because they are events on different coordinates. Linearity and Proposition~\ref{prop:varsum}(c) give $\mathbb ES=np$ and $\operatorname{Var}S=np(1-p)$.

(c) Take the $\binom Nn$ subsets of size $n$ as equally likely. Those with $k$ marked objects are chosen by picking $k$ marked objects and $n-k$ unmarked ones. For the moments, draw the objects one at a time, a uniformly random ordered selection, which induces the uniform law on subsets. Let $I_i$ be the indicator that the $i$th draw is marked. The $i$th object drawn is uniform among the $N$, by symmetry among the positions, so $P(I_i=1)=K/N$. For $i\ne j$, the pair of objects drawn at times $i$ and $j$ is uniform among the $N(N-1)$ ordered pairs, so $P(I_i=I_j=1)=K(K-1)/(N(N-1))$. Then $\mathbb EX=nK/N$, and
\[
 \mathbb E[X^2]=\sum_i\mathbb E[I_i]+\sum_{i\ne j}\mathbb E[I_i I_j]
 =n\frac KN+n(n-1)\frac{K(K-1)}{N(N-1)}.
\]
With $q=K/N$ this gives
\[
\begin{aligned}
\operatorname{Var}X
&=nq\Bigl(1+\frac{(n-1)(K-1)}{N-1}-nq\Bigr)\\
&=\frac{nq}{N-1}\bigl(N-1+nK-n-K+1-nK+nq\bigr)\\
&=\frac{nq}{N-1}(N-n)(1-q),
\end{aligned}
\]
where the middle step uses $nq(N-1)=nK-nq$.

(d) $\mathbb EX=\frac1n\sum_{k=1}^n k=(n+1)/2$, and $\mathbb E[X^2]=\frac1n\sum_{k=1}^n k^2=(n+1)(2n+1)/6$ by Chapter~\ref{ch:sums}. Subtract the square of the mean.
\end{proof}

The hypergeometric mass is zero unless $\max(0,n-N+K)\le k\le\min(n,K)$, because both chosen groups must have admissible sizes. For the excluded case $N=1$, the draw is deterministic, whether $n=0$ or $n=1$, so the variance is $0$. The formula with denominator $N-1$ is not used. The factor $(N-n)/(N-1)$ in (c) is the \emph{finite-population correction}. For $1\le n\le N$, drawing without replacement is no more variable than drawing with replacement. For $n=0$ both variances are zero. For $n=N$, drawing everything, the variance is $0$, as it must be. When $N$ is huge beside $n$, the correction is close to $1$. That is why a sample of $1000$ people is about as informative in a country of $60$ million as in one of $300$ million. Exercise~\ref{ex:poll} makes the comparison a bound.

\begin{figure}[ht]
\centering
\fig{ch08-binomial}
\caption{Left: the law of the frequency $S/n$ of heads in $n$ tosses of a fair coin, for $n=10$, $50$ and $200$. The bars have width $1/n$ and are scaled so that areas can be compared. The mass concentrates around $1/2$, in a window of width about $1/\sqrt n$. The standard deviation of $S/n$ is $1/(2\sqrt n)$. Right: the running frequency of heads in three simulated sequences of $1000$ tosses, with the band $1/2\pm 1/(2\sqrt n)$ in grey. The paths illustrate concentration at a fixed sample size. A simulation does not prove that an infinite path converges. Theorem~\ref{thm:wlln} bounds the error probability at each fixed $n$.}
\label{fig:binomial}
\end{figure}

\begin{think}
A fair coin is tossed $100$ times. You know the mean $50$ and the standard deviation $5$, and nothing else about the shape. How surprised should you be to see $70$ or more heads, or $30$ or fewer? The next theorem turns those two numbers into a bound.
\end{think}

\begin{theorem}[Markov and Chebyshev]\label{thm:chebyshev}
Let $X$ be a random variable and let $a>0$.
\begin{enumerate}[label=(\alph*)]
\item (\textbf{Markov}) If $X\ge0$, then $P(X\ge a)\le \mathbb EX/a$.
\item (\textbf{Chebyshev}) $P\bigl(|X-\mathbb EX|\ge a\bigr)\le \operatorname{Var}X/a^2$.
\end{enumerate}
\end{theorem}
\begin{proof}
(a) The inequality $a\,\mathbf1_{X\ge a}\le X$ holds at every outcome. If $X(\omega)\ge a$, the left side is $a$, and $a\le X(\omega)$. Otherwise the left side is $0$, and $0\le X(\omega)$. Take expectations. They preserve inequalities, because a sum of nonnegative terms is nonnegative. Then use $\mathbb E[\mathbf1_A]=P(A)$.

(b) Apply (a) to the nonnegative variable $(X-\mathbb EX)^2$ and the threshold $a^2$. The event $(X-\mathbb EX)^2\ge a^2$ is the event $|X-\mathbb EX|\ge a$.
\end{proof}

For the coin, $S$ has mean $50$ and variance $25$. Chebyshev with $a=20$ gives
\[
 P(|S-50|\ge20)\le\frac{25}{400}=\frac1{16}.
\]
Adding the binomial masses for $k\le30$ and for $k\ge70$ gives about $7.85\times10^{-5}$, roughly eight hundred times smaller than $1/16$. The inequality is crude for this particular law. Its use is that it holds for every random variable with a variance, however strange the law, and that it is sharp in that generality. Exercise~\ref{ex:poll}(c) exhibits the equality case. The central limit theorem in Volume~4 explains why this coin's tail is so much smaller. That explanation is not a theorem of this chapter.

\begin{theorem}[weak law of large numbers]\label{thm:wlln}
Let $n\ge1$, and let $X_1,\ldots,X_n$ be pairwise independent random variables with the same expectation $m$ and the same variance $\sigma^2$. Write $\overline X_n=(X_1+\cdots+X_n)/n$. Then for every $\varepsilon>0$,
\[
 P\bigl(|\overline X_n-m|\ge\varepsilon\bigr)\le\frac{\sigma^2}{n\varepsilon^2}.
\]
In particular (\textbf{Bernoulli's theorem}), the frequency $S/n$ of successes in $n$ independent trials of probability $p$ satisfies
\[
 P\Bigl(\Bigl|\frac Sn-p\Bigr|\ge\varepsilon\Bigr)
 \le\frac{p(1-p)}{n\varepsilon^2}
 \le\frac1{4n\varepsilon^2}.
\]
The right-hand side tends to $0$ as $n\to\infty$.
\end{theorem}
\begin{proof}
Linearity gives $\mathbb E\overline X_n=m$. Proposition~\ref{prop:varsum}(b) and (c) give $\operatorname{Var}\overline X_n=\sigma^2/n$. Apply Chebyshev. For the trials, $\sigma^2=p(1-p)\le1/4$ by Proposition~\ref{prop:laws}(a).
\end{proof}

Jacob Bernoulli spent about twenty years on this theorem. His own argument estimates binomial coefficients directly, and it fills many pages. Chebyshev's inequality reduces the proof to the four lines above, and pays for the reduction with a worse bound.

The theorem is what gives a probability a check in the world. The sentence ``this coin has probability $1/2$ of heads'' can be compared with a long series of tosses, and the theorem says how the comparison fails. Read what it does not say. It does not say that the number of heads gets close to $n/2$. The typical gap $|S-n/2|$ grows like $\sqrt n$, as in Figure~\ref{fig:binomial}. It does not say that a run of tails makes heads due. The trials are independent, and the coin has no memory of them. It says that the \emph{proportion} concentrates. The probability of a fixed proportional error tends to zero as the sample size grows. That mode of convergence is called \emph{convergence in probability}.

A strong law says something else: that the proportion converges for almost every infinite sequence. Volume~4 proves one, under a fourth-moment hypothesis, and proves the central limit theorem. This chapter compares a sequence of finite experiments, one for each sample size.

\begin{reminder}[what a sampling bound promises]
For independent Bernoulli trials, Chebyshev guarantees an error probability at most $1/(4n\varepsilon^2)$ at the chosen size $n$. For tolerance $0.05$ and failure probability at most $0.05$, the bound $n\ge2000$ suffices. The bound does not correct a biased sample, a success rate that changes, dependence, or a habit of repeating the experiment until a convenient answer appears. A bound greater than $1$ is true and useless. Replace it by the minimum of that bound and $1$.
\end{reminder}

\portraits{chebyshev}{Pafnuty Chebyshev (1821--1894). Photograph, c.\,1890. Public domain, via Wikimedia Commons.}{erdos}{Paul Erd\H os (1913--1996) in Budapest, 1992. Photograph by Topsy Kretts, CC BY 3.0, via Wikimedia Commons.}

\section{Waiting times}\label{sec:countable-probability}

No finite sample space can hold a waiting time that is allowed to be arbitrarily large. The extension used below is a convergent series of masses. It is not a theory of arbitrary infinite experiments.

A set is \emph{countable} if it is finite, or if its elements can be listed in a sequence without repetitions. On $\Omega=\{\omega_1,\omega_2,\ldots\}$, take masses $p_j\ge0$ with $\sum_j p_j=1$, and define $P(A)=\sum_{\omega_j\in A}p_j$. For a nonnegative variable $X$, put
\[
 \mathbb EX=\sum_j X(\omega_j)\,p_j,
\]
and allow the value $+\infty$. A nonnegative sum is the supremum of its finite sub-sums. That is why its terms may be grouped or reordered: every finite sub-sum sits inside a large enough partial sum in the other order. For a nonnegative double array, both iterated sums equal the supremum of the finite rectangle sums. These are the nonnegative-sum principles from Chapter~\ref{ch:sequences}. You may not cancel a divergent positive part against a divergent negative part. Finite linearity and monotonicity extend to these nonnegative expectations by the same finite-sub-sum argument. For a real variable of either sign, require $\sum_j |X(\omega_j)|p_j<\infty$ before you define the expectation by the signed series. Absolute convergence then gives linearity, by taking limits of finite partial sums. Whenever a formula below subtracts, the moments satisfy that requirement.

For a variable $T$ with values in $\{1,2,\ldots\}$, the finite identity $\min(T,M)=\sum_{j=0}^{M-1}\mathbf1_{T>j}$ gives, on taking expectations,
\[
 \mathbb E\min(T,M)=\sum_{j=0}^{M-1}P(T>j).
\]
As $M$ grows, these expectations tend to $\mathbb ET$. They are bounded above by it, and any finite sub-sum of $\sum_t t\,P(T=t)$ is captured once $M$ exceeds its indices. Hence the \emph{tail-sum identity}
\begin{equation}
 \mathbb ET=\sum_{j\ge0}P(T>j) \tag{*}
\end{equation}
holds, including when both sides are infinite.

\begin{examplex}[geometric waiting]
For $0<p\le1$, a first success on independent trials of probability $p$ has masses $P(T=t)=p(1-p)^{t-1}$ for integers $t\ge1$. For $p=1$, this says $T=1$ with certainty. For $0<p<1$, the masses add to one by the geometric series. The survival probabilities are $P(T>j)=(1-p)^j$, and the case $j=0$ is $1$. The tail-sum identity gives
\[
 \mathbb ET=\sum_{j\ge0}(1-p)^j=\frac1p.
\]
Figure~\ref{fig:geomwait} draws the masses for $p=1/2$, whose mean is $2$. This is a countable model built from its masses. No infinite sample space of coin sequences has been assumed.
\end{examplex}

\begin{figure}[ht]
\centering
\fig{ch09-geometric}
\caption{The masses of a geometric waiting time with success probability $1/2$. The bar at $t$ has height $2^{-t}$. The label $2$ is set in rust because the mean is $2$. The tail-sum identity adds the survival probabilities $P(T>j)=2^{-j}$ and gets $2$, which is $1/p$.}
\label{fig:geomwait}
\end{figure}

Check the mean another way, for $p=1/2$. With probability $1/2$ the first trial succeeds and $T=1$. With probability $1/2$ it fails, and the remaining wait is a fresh copy of $T$. So $\mathbb ET=\tfrac12\cdot1+\tfrac12\cdot(1+\mathbb ET)$, which rearranges to $\mathbb ET=2$. The fresh-copy sentence is the same geometric series in disguise. It is legal here because the tail sum has already shown that the expectation is finite.

\begin{examplex}[the Poisson law]
For $\lambda>0$, the masses $P(X=k)=e^{-\lambda}\lambda^k/k!$, for integers $k\ge0$, add to one by the exponential series. The law is called \emph{Poisson with parameter $\lambda$}. Shifting the index in the nonnegative sums gives $\mathbb EX=\lambda$ and $\mathbb E[X(X-1)]=\lambda^2$, hence $\operatorname{Var}X=\lambda$. Explicitly,
\[
 \sum_{k\ge1}k\,e^{-\lambda}\frac{\lambda^k}{k!}
 =\lambda e^{-\lambda}\sum_{j\ge0}\frac{\lambda^j}{j!}=\lambda.
\]
The second factorial moment follows by shifting twice, and $X^2=X(X-1)+X$. For $\lambda=0$ the law is concentrated at $0$. The bars in Figure~\ref{fig:derang} compare these masses, at $\lambda=1$, with the fixed-point law of a random permutation. The agreement on the page is numerical. The proof that the laws converge belongs to Volume~4.
\end{examplex}

The remaining arguments of the exercises stay finite, except where a countable model is explicitly built. The weak law compares a sequence of finite experiments, one for each sample size. One probability space that supports an infinite sequence, and almost-sure statements on it, belongs to Volume~4.

\begin{beauty}
The exact count $D_n/n!$ is Proposition~\ref{prop:derang}, from the counting chapter. The Poisson masses are the example above. What the indicator sum buys is a proof that does not build the law. Exercises~\ref{ex:fixedvar}, \ref{ex:records} and~\ref{ex:coupon} are three instances. Exercise~\ref{ex:ramsey} is Erd\H os's 1947 use of the same idea: if the expected number of monochromatic cliques is less than $1$, then some colouring has none. The argument shows that a colouring exists. It does not exhibit one. An explicit colouring with the same exponential lower bound is still missing. The challenge under the exercises records how little that lower bound has moved. Chebyshev's inequality is the other purchase. A poll, a clinical trial, and a repeated experiment are Theorem~\ref{thm:wlln} with a chosen tolerance and a chosen failure probability. The reminder above is the budget.
\end{beauty}

\section{What has been proved}

A probability, in this chapter, is a weighted count on a set you chose before you computed. A conditional probability is a new weighting of the same outcomes, not a new logic. Expectation is linear whether or not the variables are independent. Variance adds when the summands are pairwise uncorrelated, and Chebyshev turns a variance into a numerical guarantee. The weak law is that guarantee applied to an average. The hats, the shared birthday, and the three doors are one method. Write the experiment, then the weights, then the question.

Put the two sentences of the book side by side. The sum of cubes equals the square of the triangular number for every $n$. That is an identity. Both sides name the same integer. The hat check does not. For $n=2$ the probability of a derangement is $1/2$. For $n=10$ it already lies within $3\times10^{-8}$ of $1/e$. The limit is a statement about a sequence of finite models. It is not a claim that some infinite party has probability exactly $1/e$. You can now say which sentence you have proved, because you can point at the quantifiers.

A uniformly chosen integer from $1$ to $N$ is a finite probability model, even though the integer, once chosen, is a definite number. For each fixed prime $p$, the proportion of integers up to $N$ that are divisible by $p$ is $\lfloor N/p\rfloor/N$, which tends to $1/p$. Exercise~\ref{ex:arithmetic-independence} lets several distinct primes act at once: exactly, when $N$ is a multiple of their product, and with an explicit error when it is not. The argument uses the unique factorisation from Chapter~\ref{ch:sums}, inclusion--exclusion, and a limit. It keeps the set of primes fixed.

That is as far as this book goes. Volume~2, \emph{The Integers: Primes, Patterns and Proofs}, asks how often an integer up to $N$ is itself prime, and what the divisibility in that exercise becomes when the primes are the object rather than a fixed list. This book does not answer that question. Volume~4, \emph{Chance: The Logic of Uncertainty}, builds the limit theorems that need an infinite sequence of trials, including a strong law under a fourth-moment hypothesis and the central limit theorem. The field guide after this chapter gathers the habits of proof. The primes are the next subject.

\section{Exercises}

\begin{exercise}[$\star$]\label{ex:birthday}
\emph{The birthday problem.} Show that the probability that $n$ people (with birthdays uniform and independent among 365 days) all have different birthdays is at most
$e^{-n(n-1)/730}$, and deduce that it is less than $\frac12$ for $n=23$.
\end{exercise}

\begin{exercise}[$\star\star$]\label{ex:fixedvar}
Let $X$ be the number of fixed points of a uniformly random permutation of $\{1,\dots,n\}$, $n\ge2$. Show that $\operatorname{Var}X=1$.
\end{exercise}

\begin{exercise}[$\star\star$]\label{ex:cycles}
Let $n\ge1$ and let $C$ be the number of cycles of a uniformly random permutation of $\{1,\dots,n\}$.
\begin{enumerate}[label=(\alph*)]
\item Show that the length $L$ of the cycle containing $1$ is uniform on $\{1,\dots,n\}$.
\item Show that $C=\sum_{i=1}^n\frac1{L_i}$, where $L_i$ is the length of the cycle containing $i$, and deduce $\mathbb E[C]=H_n\sim\ln n$.
\end{enumerate}
\end{exercise}

\begin{exercise}[$\star\star\star$]\label{ex:coupon}
\emph{The coupon collector.} Let $n\ge1$. Each cereal box contains one of $n$ coupons, uniformly at random and independently. Show that the expected number of boxes needed to collect all $n$
coupons is $nH_n\sim n\ln n$. \emph{Use the geometric waiting law of Section~\ref{sec:countable-probability}. The solution constructs the complete stopping-time law from finite words, so no probability space of infinite trials is needed.}
\end{exercise}

\subsection*{Further exercises}

\begin{exercise}[$\star$]\label{ex:monty}
\emph{Monty Hall, with variations.} The car is placed uniformly behind the three doors. You pick door $1$ and observe the host open door $3$, revealing a goat. In part (b), $0\le q\le1$ and the informed host always reveals a goat.
\begin{enumerate}[label=(\alph*)]
\item \emph{The forgetful host} does not know where the car is and opens one of the two other doors at random; it happened to show a goat. Show that switching now wins with probability $\frac12$.
\item \emph{The biased host} knows where the car is and, when he has a choice, opens door $3$ with probability $q$ (and door $2$ with probability $1-q$). Show that switching wins with probability $\frac1{1+q}$, and that it is never worse than staying.
\item With $100$ doors, assume the car is placed uniformly at random and the informed host always opens $98$ goat doors. What is the unconditional probability of winning with the strategy of always switching?
\end{enumerate}
\end{exercise}

\begin{exercise}[$\star\star$]\label{ex:multinomial}
\emph{Bridge: Counting, Stirling and probability.}
\emph{Multinomial coefficients.} Let $r\ge1$, $n\ge0$ and $k_1,\ldots,k_r\ge0$ be integers with $k_1+\dots+k_r=n$.
\begin{enumerate}[label=(\alph*)]
\item Show that the number of ways to split $\{1,\dots,n\}$ into labelled boxes of sizes $k_1,\dots,k_r$ is $\binom n{k_1,\dots,k_r}=\frac{n!}{k_1!\cdots k_r!}$, and that this is also the number of words with $k_i$ letters of type $i$. How many anagrams has MISSISSIPPI?
\item Prove the multinomial theorem $(x_1+\dots+x_r)^n=\sum\binom n{k_1,\dots,k_r}x_1^{k_1}\cdots x_r^{k_r}$, the sum over all $(k_i)$ with $\sum k_i=n$, and deduce the number of terms.
\item In $n$ independent rolls of a fair die, what is the probability of seeing each face exactly $n/6$ times (for $n$ a multiple of $6$)? Estimate it for $n=600$ with Stirling's formula.
\end{enumerate}
\end{exercise}

\begin{exercise}[$\star\star$]\label{ex:bonferroni}
\emph{Bonferroni's inequalities.} For $n\ge1$ events $A_1,\dots,A_n$, show that
\[
\sum_iP(A_i)-\sum_{i<j}P(A_i\cap A_j)\ \le\ P\Big(\bigcup_iA_i\Big)\ \le\ \sum_iP(A_i).
\]
(Prove the corresponding inequalities between indicator functions at a point lying in exactly $m$ of the sets.) Use the upper bound, the \emph{union bound}, to show that if $n$ independent events each have probability $\le\frac1{2n}$, then with probability at least $\frac12$ none of them occurs.
\end{exercise}

\begin{exercise}[$\star\star$]\label{ex:records}
\emph{Records.} For $n\ge1$, in a uniformly random permutation $(a_1,\dots,a_n)$ of $\{1,\dots,n\}$, position $i$ is a \emph{record} if $a_i>a_j$ for all $j<i$. Let $R_i$ be the indicator of a record at $i$ and $R=\sum R_i$ the number of records.
\begin{enumerate}[label=(\alph*)]
\item Show that $P(R_i=1)=\frac1i$, and deduce $\mathbb ER=H_n$.
\item Show that for $i<j$, $P(R_i=R_j=1)=\frac1{ij}$: the record indicators are pairwise independent. Deduce $\operatorname{Var}R=H_n-\sum_{i\le n}\frac1{i^2}$.
\item Use Chebyshev to show that the number of records among $10^6$ random values is between $4$ and $25$ with probability at least $0.9$.
\end{enumerate}
\end{exercise}

\begin{exercise}[$\star\star$]\label{ex:poll}
\emph{A poll.} A fraction $p$ of a large population intends to vote for a candidate. You ask $n$ people chosen at random (with replacement, say) and record the fraction $\overline X_n$ who say yes.
\begin{enumerate}[label=(\alph*)]
\item Show that $P(|\overline X_n-p|\ge0.03)\le\frac1{4n(0.03)^2}$, and find $n$ such that this is at most $0.05$.
\item Why does the size of the population not appear? What changes if you sample without replacement?
\item Show that Chebyshev's inequality cannot be improved in general: for every $a>0$ exhibit a random variable with $P(|X-\mathbb EX|\ge a)=\frac{\operatorname{Var}X}{a^2}$.
\end{enumerate}
\end{exercise}

\begin{exercise}[$\star\star\star$]\label{ex:ballot}
\emph{Bertrand's ballot problem.} Let $a>b\ge0$ be integers. In an election, $A$ receives $a$ votes and $B$ receives $b$ votes. The ballots are counted in uniformly random order. Show that the probability that $A$ is strictly ahead throughout the count is $\frac{a-b}{a+b}$. (Represent a count as a path with steps $+1$ for $A$ and $-1$ for $B$. Count the bad paths that start with a $+1$ by reflecting the part of the path up to its first return to $0$: the reflection argument often associated with Andr\'e's 1887 solution, although his original bijection was different; see Renault in the sources.)
\end{exercise}

\begin{exercise}[$\star\star\star$]\label{ex:ramsey}
\emph{Erd\H os's lower bound for Ramsey numbers.} Let $n\ge1$ and $k\ge2$ be integers, with the count in (a) understood as zero when $k>n$. The \emph{complete graph} has one edge for each unordered pair of distinct vertices. Colour each of the $\binom n2$ edges of the complete graph on $n$ vertices red or blue, independently with probability $\frac12$ each.
\begin{enumerate}[label=(\alph*)]
\item Show that the expected number of sets of $k$ vertices whose $\binom k2$ edges all have the same colour is $\binom nk2^{1-\binom k2}$.
\item Show that this is $<1$ when $n=\lfloor2^{k/2}\rfloor$ and $k\ge3$, and conclude that there is a colouring of the complete graph on $n$ vertices with no monochromatic set of $k$ vertices: in the notation of Ramsey theory, $R(k,k)>2^{k/2}$.
\item Why does the argument prove existence without producing a colouring? For $k=3$, find by hand a colouring of the complete graph on $5$ vertices with no monochromatic triangle, and show that every colouring on $6$ vertices has one: $R(3,3)=6$.
\end{enumerate}
\end{exercise}

\begin{exercise}[$\star\star$]\label{ex:covariance-sampling}
\emph{Bridge: counting, eigenvalues and dependence.} Choose a $k$-element subset
of $\{1,\ldots,n\}$ uniformly, where $n\ge2$ and $0\le k\le n$, and let
$X_i$ indicate that $i$ is chosen. Set $p=k/n$ and let $J$ be the $n\times n$
matrix all of whose entries are one.
\begin{enumerate}[label=(\alph*)]
\item Find $\mathbb EX_i$ and $\operatorname{Cov}(X_i,X_j)$, including $i=j$.
\item Prove that the covariance matrix is
$\Sigma=\dfrac{p(1-p)}{n-1}(nI-J)$. Find all its eigenvalues and its rank.
\item Explain the zero eigenvalue directly from the experiment. Why are distinct
indicators negatively correlated when $0<k<n$?
\end{enumerate}
\end{exercise}

\begin{exercise}[$\star\star$]\label{ex:arithmetic-independence}
\emph{Bridge: divisibility, inclusion--exclusion and chance.} Let
$p_1,\ldots,p_r$ be distinct primes, $r\ge1$, and $q=p_1\cdots p_r$.
Choose $U$ uniformly from $\{1,\ldots,N\}$, where $N\ge1$.
\begin{enumerate}[label=(\alph*)]
\item If $q$ divides $N$, prove that the events $p_i\mid U$ are mutually
independent, with probabilities $1/p_i$. Find $P(\gcd(U,q)=1)$.
\item For general $N$, prove
\[
 \left|P(\gcd(U,q)=1)-\prod_{i=1}^r(1-1/p_i)\right|
 \le\frac{2^r-1}{N}.
\]
\item Give an example in which $q$ does not divide $N$ and the divisibility
events are not independent. What does part (b) say as $N\to\infty$ with the
primes fixed?
\end{enumerate}
\end{exercise}

\begin{challenge}
\textbf{Three questions to take further.} (1) \emph{The problem of the m\'enages} (Lucas, 1891): for $n\ge3$, first fix the cyclic order and positions of $n$ women around a table. In how many ways can their $n$ husbands fill the intervening gaps, no one next to their spouse? Touchard's 1934 formula for this number is $U_n=\sum_{k=0}^n(-1)^k\frac{2n}{2n-k}\binom{2n-k}k(n-k)!$. If the women's cyclic order is also free and rotations count as the same seating, the total is $(n-1)!U_n$; mirror images remain distinct. (2) \emph{Ramsey numbers.} $R(3,3)=6$ is a party puzzle (among six people, three know each other or three are mutual strangers); $R(4,4)=18$ was found in 1955; $R(5,5)$ is unknown, between $43$ and $46$ as of October 2026. Erd\H os liked to say that if aliens demanded $R(5,5)$ on pain of destroying the Earth, we should put every computer to work, but if they asked for $R(6,6)$, we should attack. Why is the exponential lower bound of Exercise~\ref{ex:ramsey}, found in 1947, still essentially the best known? (3) \emph{Expected waiting for a shared birthday.} Show that the expected number of people one must gather until two share a birthday is $\sum_{k=0}^{365}\prod_{j=0}^{k-1}\big(1-\frac j{365}\big)\approx24.6$. For a year of $N$ days, the advanced asymptotic refinement is $\sqrt{\pi N/2}+\frac23+o(1)$. This refinement is quoted as an outlook, not proved by the heuristic in the notes; a uniform error analysis is needed in addition to the calculus of Chapter~\ref{ch:calculus}.
\end{challenge}

\begin{exercise}[$\star\star$]\label{pb:cauchydav}
\emph{Cauchy--Davenport.} Let $p$ be prime and $A,B\subset\Z/p\Z$ non-empty.
\begin{enumerate}[label=(\alph*)]
\item (Cauchy 1813, Davenport 1935.) Prove that $|A+B|\ge\min(p,|A|+|B|-1)$, using Problem~\ref{pb:nullstellensatz}.
\item Deduce $|A_1+\dots+A_k|\ge\min\bigl(p,\sum|A_i|-k+1\bigr)$.
\item Let $a_1,\dots,a_{p-1}$ be non-zero residues, not necessarily distinct. Show that every residue is the sum of some of the
$a_i$. Is $p-1$ the least number that works?
\end{enumerate}
\end{exercise}

\begin{exercise}[$\star\star$]\label{pb:distances}
\emph{Few distances, few unit distances.} Let $P$ be a set of $n\ge2$ points in the plane.
\begin{enumerate}[label=(\alph*)]
\item (Erdős, 1946.) Show that $P$ determines at least $\sqrt{(n-2)/2}$ distinct distances.
\item Show that at most $\frac n4\bigl(1+\sqrt{8n-7}\bigr)$ pairs of points of $P$ are at distance exactly $1$.
\end{enumerate}
\end{exercise}

\begin{exercise}[$\star\star$]\label{pb:incidences}
\emph{Incidences by Cauchy--Schwarz.} Let $P$ be a set of $n$ points and $L$ a set of $m$ lines in the plane, and let $I(P,L)$ be the
number of pairs $(p,\ell)$ with $p\in\ell$.
\begin{enumerate}[label=(\alph*)]
\item Show that $I(P,L)\le n\sqrt m+m$ and $I(P,L)\le m\sqrt n+n$.
\item Let $k\ge2\sqrt n$. Show that at most $\frac{2n}k$ lines contain $k$ or more points of $P$.
\end{enumerate}
\end{exercise}

\begin{exercise}[$\star\star\star$]\label{pb:szemereditrotter}
\emph{The Szemerédi--Trotter theorem (Székely's proof).} Recall the crossing lemma, proved in Volume~4: a simple graph with
$v$ vertices and $e\ge4v$ edges cannot be drawn in the plane with fewer than $\frac{e^3}{64v^2}$ crossings.
\begin{enumerate}[label=(\alph*)]
\item With the notation of Problem~\ref{pb:incidences}, show that $I(P,L)\le4(nm)^{2/3}+m+4n$. (Join consecutive points of $P$ on
each line.)
\item Deduce that, for $k\ge2$, the number of lines containing at least $k$ points of $P$ is $O\bigl(\frac{n^2}{k^3}+\frac nk\bigr)$.
\item By the same method with unit circles, show that $n$ points determine at most $4n^{4/3}+5n$ unit distances.
\end{enumerate}
\end{exercise}

\begin{exercise}[$\star\star$]\label{pb:sumproduct}
\emph{Sums or products (Elekes, 1997).} Let $A$ be a finite set of $N$ positive real numbers. By looking at the point set
$(A+A)\times(A\cdot A)$ and the $N^2$ lines $y=b(x-c)$, $b,c\in A$, show that
\[
|A+A|\cdot|A\cdot A|\ge c\,N^{5/2}
\]
for an absolute constant $c>0$, so that $\max(|A+A|,|A\cdot A|)\ge\sqrt c\,N^{5/4}$.
\end{exercise}

\begin{exercise}[$\star$]\label{pb:treebound}
\emph{Spanning trees and AM--GM.} Let $G$ be a connected graph with $n$ vertices and $m$ edges, and $\tau(G)$ its number of spanning
trees. Show that
\[
\tau(G)\le\frac1n\Bigl(\frac{2m}{n-1}\Bigr)^{n-1},
\]
with equality if and only if $G$ is complete. Which theorem do you recover for $K_n$?
\end{exercise}

\begin{exercise}[$\star\star$]\label{pb:parking}
\emph{Parking on a one-way street.} Spots $1,\dots,m$ line a one-way street. Cars arrive one by one; car $i$ prefers spot $a_i$, drives to it, and
parks in the first free spot from $a_i$ on, leaving the street if there is none. A \emph{parking function} of length $n$ is a sequence
$(a_1,\dots,a_n)\in\{1,\dots,n\}^n$ with which, for $m=n$, every car parks.
\begin{enumerate}[label=(\alph*)]
\item Put $n+1$ spots on a circle instead, with preferences in $\Z/(n+1)\Z$, so that every car parks and one spot stays empty. By rotating preferences,
show that there are $(n+1)^{n-1}$ parking functions of length $n$.
\item Back on the street, take $m=n+1$ and preferences uniform in $\{1,\dots,n+1\}^n$, and let $K+1$ be the first spot left empty. Show that
\[
P(K=k)=\binom nk\frac{(k+1)^{k-1}(n-k)^{n-k}}{(n+1)^n},
\]
and deduce Abel's identity $\sum_{k=0}^n\binom nk(k+1)^{k-1}(n-k)^{n-k}=(n+1)^n$.
\item Show that $P(K=k)\to p_k=\frac{(k+1)^{k-1}}{k!}e^{-(k+1)}$ for each fixed $k$, and that $\sum_kp_k=1$. (Use the tree function $T(z)=\sum_{m\ge1}m^{m-1}\frac{z^m}{m!}$,
which satisfies $T(z)e^{-T(z)}=z$.)
\item Show that $p_k\sim\frac1{\sqrt{2\pi}}k^{-3/2}$. So $K$ is usually small, yet the limit law has infinite mean.
\end{enumerate}
\end{exercise}

\begin{exercise}[$\star\star\star$]\label{pb:sinesecond}
\emph{The sine iteration, second term.} Exercise~\ref{ex:sin} showed that $u_{n+1}=\sin u_n$, with $0<u_0\le\frac\pi2$, satisfies $u_n\sim\sqrt{3/n}$.
\begin{enumerate}[label=(\alph*)]
\item Show that $\frac1{\sin^2u}=\frac1{u^2}+\frac13+\frac{u^2}{15}+O(u^4)$ as $u\to0$.
\item Show that $\frac1{u_n^2}=\frac n3+O(\ln n)$, then that $\frac1{u_n^2}=\frac n3+\frac15\ln n+C+o(1)$ for a constant $C$ depending on $u_0$.
\item Deduce that $u_n=\sqrt{\frac3n}-\frac{3\sqrt3}{10}\cdot\frac{\ln n}{n^{3/2}}+O\bigl(\frac1{n^{3/2}}\bigr)$.
\end{enumerate}
\end{exercise}

\section{Hints}

\begin{hint}{ex:birthday}
The probability is $\prod_{k<n}\left(1-\frac k{365}\right)$. Use $1-x\le e^{-x}$ to turn the product into the exponential of a sum.
\end{hint}

\begin{hint}{ex:fixedvar}
$X=\sum\mathbf 1_{\sigma(i)=i}$, so $X^2=\sum_{i,j}\mathbf 1_{\sigma(i)=i,\,\sigma(j)=j}$. What is the probability that both $i\ne j$ are fixed?
\end{hint}
\begin{hint}{ex:cycles}
(a) Count the permutations whose cycle through $1$ has length $\ell$: choose the $\ell-1$ other elements of the cycle in order, then permute the rest. (b) A cycle of length $\ell$ contributes $\ell$ times $\frac1\ell$ to $\sum\frac1{L_i}$.
\end{hint}
\begin{hint}{ex:coupon}
Cut the wait into phases: with $k$ distinct coupons in hand, a new box brings a new coupon with probability $\frac{n-k}n$.
\end{hint}

\begin{hint}{ex:monty}
Redo the computation of the example with the new values of $P(H_3\mid C_i)$. In (a), $P(H_3\mid C_3)$ is no longer $0$, but the event ``door $3$ shows a goat'' is.
\end{hint}
\begin{hint}{ex:multinomial}
(a) Choose the first box, then the second among what remains. (b) Expand the product and count the ways to obtain each monomial. (c) The probability is $\binom n{n/6,\dots,n/6}6^{-n}$; Stirling (Theorem~\ref{thm:stirling}).
\end{hint}

\begin{hint}{ex:bonferroni}
At a point in exactly $m\ge1$ of the sets, the three expressions are $m-\binom m2$, $1$ and $m$. For the application, use $P(A_i^c\cap\dots)=1-P(\bigcup A_i)$; independence is not even needed.
\end{hint}
\begin{hint}{ex:records}
(a) Position $i$ is a record iff the largest of $a_1,\dots,a_i$ is at position $i$; by symmetry each position is equally likely. (b) Condition on the record at $j$: the relative order of $a_1,\dots,a_{j-1}$ is still uniform. (c) Bound $H_{20}$ and $\sum_{i\le20}1/i^2$ exactly, then compare the tail with integrals. An upper bound on $\sum 1/i^2$ lowers the variance and is the wrong direction for Chebyshev.
\end{hint}
\begin{hint}{ex:poll}
(a) Theorem~\ref{thm:wlln}. (b) Proposition~\ref{prop:laws}(c). (c) Three values: $-a$, $0$, $a$.
\end{hint}
\begin{hint}{ex:ballot}
There are $\binom{a+b}a$ equally likely paths. Those that start with $-1$ are bad; show that the bad paths starting with $+1$ are in bijection with the paths starting with $-1$. Good paths: $\binom{a+b-1}{a-1}-\binom{a+b-1}{a}$.
\end{hint}
\begin{hint}{ex:ramsey}
(a) Indicators. (b) $\binom nk\le\frac{n^k}{k!}$. (c) An expectation below $1$ for a non-negative integer variable forces a zero somewhere.
\end{hint}

\begin{hint}{ex:covariance-sampling}
Count subsets containing one prescribed element, then two. The vector of all
ones and the subspace whose coordinates sum to zero are invariant under $J$.
Treat $k=0,n$ separately before dividing by a variance.
\end{hint}

\begin{hint}{ex:arithmetic-independence}
A product of distinct primes divides $U$ exactly when each of those primes does.
There are $\lfloor N/d\rfloor$ multiples of $d$ up to $N$. Apply
inclusion--exclusion before estimating each floor error.
\end{hint}

\begin{hint}{pb:cauchydav}
(a) If $A+B$ were inside a set $C$ of size $|A|+|B|-2$, the polynomial $\prod_{c\in C}(x+y-c)$ would vanish on $A\times B$.
\end{hint}

\begin{hint}{pb:distances}
(a) Fix two points; every other point is on a circle around each. (b) The unit-distance graph has no $K_{2,3}$: count pairs of
neighbours.
\end{hint}

\begin{hint}{pb:incidences}
(a) Count pairs of points on a common line, then use $\bigl(\sum_\ell k_\ell\bigr)^2\le m\sum_\ell k_\ell^2$.
\end{hint}

\begin{hint}{pb:szemereditrotter}
(a) The graph has $I-m$ edges (at least, once lines with no points are removed) and at most $m^2$ crossings. (c) Two unit circles cross
at most twice; ignore circles with fewer than $3$ points.
\end{hint}

\begin{hint}{pb:sumproduct}
Each line contains the $N$ points $(a+c,ab)$. Apply Problem~\ref{pb:szemereditrotter}(a).
\end{hint}

\begin{hint}{pb:treebound}
Matrix-tree theorem: $\tau(G)=\frac1n\mu_2\cdots\mu_n$ for the Laplacian eigenvalues.
\end{hint}

\begin{hint}{pb:parking}
(a) On the circle, adding $c$ to every preference rotates the final configuration by $c$. (b) If spot $k+1$ is the first empty one, the cars that preferred a spot at most $k+1$
are exactly the cars parked in $1,\dots,k$. (c) $T$ increases on $\left[0,\frac1e\right[$, and $te^{-t}$ increases from $0$ to $\frac1e$ on $[0,1]$. (d) Stirling.
\end{hint}

\begin{hint}{pb:sinesecond}
(a) $\sin^2u=u^2(1-\frac{u^2}3+\frac{2u^4}{45}+O(u^6))$. (b) First sum the crude estimate $\frac1{u_{k+1}^2}-\frac1{u_k^2}=\frac13+O(u_k^2)$, then the precise one. (c) $(\frac n3+a)^{-1/2}$.
\end{hint}

\section{Solutions}

\begin{solution}{ex:birthday}
\idea Count favourable outcomes as injections: the probability is a product of factors $1-\frac k{365}$. Products are hard to estimate, sums are easy: the hint $1-x\le e^{-x}$ turns the product into the exponential of a sum. (The hint itself is the tangent-line inequality for the convex function $e^{-x}$, Proposition~\ref{prop:convex}.)
\writeup For $n>365$ the probability of all distinct birthdays is $0$, so the bound is immediate. For $0\le n\le365$, with $\Omega$ the $365^n$ birthday assignments (uniform), the favourable ones are the injections, $\frac{365!}{(365-n)!}$ of them, so the probability is $\prod_{k=0}^{n-1}\left(1-\frac k{365}\right)\le\exp\left(-\sum_{k=0}^{n-1}\frac k{365}\right)=e^{-n(n-1)/730}$.
For $n=23$, $n(n-1)/730=253/365$, so the bound is $e^{-253/365}$. The logarithm workshop in Chapter~\ref{ch:sequences} gives $\ln2<0.693148$. The comparison $693148\cdot365=252999020<253\cdot10^6$ says $0.693148<253/365$, hence $\ln2<253/365$. The exponent $-253/365$ is then less than $\ln(1/2)$, and the bound is below $\frac12$.
The exact probability is about $0.4927$.
\end{solution}

\begin{solution}{ex:fixedvar}
\idea The indicator trick computes $\mathbb E[X]$; for the variance we need $\mathbb E[X^2]$, and squaring a sum of indicators gives indicators of \emph{pairs} of events. So we need the probability that two given points are both fixed, a simple count.
\writeup Write $X=\sum_i\mathbf 1_{\sigma(i)=i}$. Then
\[
\begin{aligned}X^2&=\sum_i\mathbf 1_{\sigma(i)=i}\eqbr+\sum_{i\ne j}\mathbf 1_{\sigma(i)=i,\,\sigma(j)=j}\end{aligned}
\]
(as $\mathbf 1_A^2=\mathbf 1_A$ and $\mathbf 1_A\mathbf 1_B=\mathbf 1_{A\cap B}$). A permutation fixing $i\ne j$ is a permutation of the other $n-2$ points,
so $P(\sigma(i)=i,\sigma(j)=j)=\frac{(n-2)!}{n!}$. Hence $\mathbb E[X^2]=1+n(n-1)\frac{(n-2)!}{n!}=2$ and $\operatorname{Var}X=2-1^2=1$.
\end{solution}

\begin{solution}{ex:cycles}
\idea (a) To show a law is uniform, count the permutations with $L=\ell$ and observe the count does not depend on $\ell$. Build such a permutation by first walking along the cycle of $1$, then permuting the rest. (b) Counting cycles directly is awkward, but each element can ``own'' a share $\frac1\ell$ of its cycle of length $\ell$, so the shares of a cycle add up to $1$. This rewrites $C$ as a sum, and linearity applies.
\writeup (a) A permutation whose cycle through $1$ has length $\ell$ is given by the ordered list $\sigma(1),\sigma^2(1),\dots,\sigma^{\ell-1}(1)$ of distinct elements other than $1$
($\frac{(n-1)!}{(n-\ell)!}$ choices) and an arbitrary permutation of the $n-\ell$ remaining points ($(n-\ell)!$ choices): $(n-1)!$ permutations in all, for each $\ell$. So $P(L=\ell)=\frac{(n-1)!}{n!}=\frac1n$.

(b) Each cycle of length $\ell$ contributes $\ell\cdot\frac1\ell=1$ to $\sum_i\frac1{L_i}$, so the sum counts cycles. By symmetry (relabel $i$ as $1$) every $L_i$ has the law of $L$, and by linearity
$\mathbb E[C]=n\,\mathbb E\left[\frac1L\right]=n\cdot\frac1n\sum_{\ell=1}^n\frac1\ell=H_n$, which is $\sim\ln n$ by Proposition~\ref{prop:gamma}.
\end{solution}

\begin{solution}{ex:coupon}
\idea The last coupons are hard to obtain. Split the wait into phases, from $k$
to $k+1$ distinct coupons: their success probabilities are
$1,(n-1)/n,\ldots,1/n$. To justify an unlimited wait using only what we have
built, count finite words with prescribed phase lengths and check that their
probabilities sum to one. Then the geometric expectation and finite linearity
finish the calculation.
\writeup Let $t_k\ge1$ be the number of boxes bought while holding exactly $k$
distinct coupons, ending with the next new coupon, for $0\le k<n$.
For prescribed lengths $(t_0,\ldots,t_{n-1})$, the finite coupon word has
$t_0+\cdots+t_{n-1}$ letters. In phase $k$, each of its first $t_k-1$
letters has $k$ choices among the old coupons, and its last letter has $n-k$
choices among the new ones. Consequently the total probability of these
words is
\[
 \prod_{k=0}^{n-1}\left(\frac{k}{n}\right)^{t_k-1}\frac{n-k}{n}.
\]
For $k=0$, this factor is one when $t_0=1$ and zero otherwise: the first box
always starts the collection. Each factor is the geometric mass of parameter
$(n-k)/n$; summing the product over all tuples gives one, by successive
nonnegative summation and the geometric series. The tuples can be listed by increasing total length. These masses therefore define
a countable probability model for the completed collection, directly from its
finite histories. For any $M$, a stopped word of length $t\le M$ has
$n^{M-t}$ extensions to $M$ draws, each of mass $n^{-M}$; summing recovers
exactly the finite experiment's probability of completing the collection by
$M$. No mass is left for a collection that never finishes.

The phase lengths are independent in this model, with
$\mathbb ET_k=n/(n-k)$ by the geometric waiting example; independence is
not needed to add their means. For $T=\sum_{k=0}^{n-1}T_k$, nonnegative
linearity gives
\[
 \mathbb ET=\sum_{k=0}^{n-1}\frac{n}{n-k}
 =n\sum_{j=1}^n\frac1j=nH_n\sim n\ln n.
\]
This includes $n=1$, when $T=1$. Equivalently, the expectation is the limit of
finite truncated waits, as established in Section~\ref{sec:countable-probability};
it is not obtained by applying finite linearity to an unjustified infinite sum
of indicators.
\end{solution}

\begin{solution}{ex:monty}
\idea Nothing changes in the method: list the hypotheses $C_1,C_2,C_3$, write the probability of the observed evidence under each, and apply Bayes. What changes is the evidence. In (a) the evidence is ``door $3$ was opened \emph{and} showed a goat'', which is no longer certain under $C_1,C_2$ and impossible under $C_3$ only because of the goat, not because of the host's knowledge.
\writeup (a) The forgetful host opens door $2$ or $3$ with probability $\frac12$ each, whatever the car does. The evidence is $G=H_3\cap C_3^c$ (door $3$ opened, goat seen). $P(G)=P(H_3)-P(H_3\cap C_3)=\frac12-\frac12\cdot\frac13=\frac13$, and $P(C_2\mid G)=\frac{P(C_2\cap H_3)}{P(G)}=\frac{\frac13\cdot\frac12}{\frac13}=\frac12$. Switching and staying are now equivalent: a host who could have revealed the car, and did not, carries only the information ``the car is not behind door $3$'', which is symmetric in doors $1$ and $2$.

(b) Now $P(H_3\mid C_1)=q$, $P(H_3\mid C_2)=1$, $P(H_3\mid C_3)=0$, so $P(H_3)=\frac q3+\frac13$ and
\[
P(C_2\mid H_3)=\frac{1\cdot\frac13}{\frac{1+q}3}=\frac1{1+q}\ \ge\ \frac12 ,
\]
with equality only for $q=1$. For a general $q$,
\[
P(C_1\mid H_3)=\frac{q\cdot\frac13}{\frac{1+q}3}=\frac q{1+q}.
\]
This equals $\frac12$ only when $q=1$. In that case $P(C_1\mid H_3)=P(C_2\mid H_3)=\frac12$, so switching is no better than staying. The observation has changed the probability of door $1$ from $\frac13$ to $\frac12$, and doors $1$ and $2$ are then tied. For $q=0$, the host opening door $3$ means he had no choice, $P(C_1\mid H_3)=0$, and switching wins surely.

(c) The car is behind door $1$ with probability $\frac1{100}$. Whatever happens, the host leaves exactly one of the other $99$ doors closed, and if the car is among those $99$ doors (probability $\frac{99}{100}$) it is behind the closed one, since he may only open goat doors. So switching wins with probability $\frac{99}{100}$: the host has done the contestant's search for them. This is the unconditional success probability of the strategy. Conditioning on \emph{which} $98$ doors he opened requires the host's selection rule; without a symmetric rule, the conditional probability need not be $\frac{99}{100}$.
\end{solution}

\begin{solution}{ex:multinomial}
\idea (a) Fill the boxes one after the other with the product rule; the binomial coefficients telescope. The word count is the same object seen as a map ``position $\mapsto$ letter''. (b) The expansion of a power is a sum over words, exactly as for the binomial theorem. (c) The probability is a single multinomial coefficient over $6^n$, and Stirling's formula turns factorials into powers.
\writeup (a) Choose the $k_1$ elements of box $1$ ($\binom n{k_1}$ ways), then the $k_2$ elements of box $2$ among the remaining $n-k_1$, and so on:
\[
\binom n{k_1}\binom{n-k_1}{k_2}\cdots\binom{k_r}{k_r}=\frac{n!}{k_1!(n-k_1)!}\cdot\frac{(n-k_1)!}{k_2!(n-k_1-k_2)!}\cdots=\frac{n!}{k_1!\cdots k_r!}.
\]
A word of length $n$ with $k_i$ letters of type $i$ is the same as the assignment of each position to the box of its letter. MISSISSIPPI has $11$ letters, with multiplicities $1,4,4,2$: $\frac{11!}{1!\,4!\,4!\,2!}=34650$ anagrams.

(b) Expanding $(x_1+\dots+x_r)^n$ chooses one $x_i$ from each of the $n$ factors: each word of length $n$ in the letters $1,\dots,r$ contributes one term, and the word with $k_i$ letters $i$ contributes $x_1^{k_1}\cdots x_r^{k_r}$. Grouping by the multiplicities gives the multinomial coefficient. The number of monomials is the number of $(k_1,\dots,k_r)\in\N^r$ with sum $n$: $\binom{n+r-1}{r-1}$ by Exercise~\ref{ex:stars}.

(c) Each of the $6^n$ sequences of rolls is equally likely, and those with each face $m=n/6$ times number $\binom n{m,\dots,m}$. By Stirling's formula (Theorem~\ref{thm:stirling}), $n!\sim\sqrt{2\pi n}(n/e)^n$ and $m!^6\sim(2\pi m)^3(m/e)^{6m}=(2\pi n/6)^3(n/e)^n6^{-n}$, so
\[
\frac{n!}{m!^6\,6^n}\sim\frac{\sqrt{2\pi n}}{(2\pi n/6)^3}=\frac{216}{(2\pi n)^{5/2}} ,
\]
about $2.5\times10^{-7}$ for $n=600$ (the exact value is $2.46\times10^{-7}$). The ``most likely'' outcome is very unlikely: the probability is spread over the roughly $\binom{605}5\approx6.6\times10^{11}$ possible count vectors, and even the peak gets only a tiny share.
\end{solution}

\begin{solution}{ex:bonferroni}
\idea Inclusion--exclusion was proved pointwise, by an identity between indicator functions; truncating that identity gives inequalities, also pointwise, and taking expectations turns them into inequalities between probabilities. The whole content is the behaviour of $m-\binom m2$ for integers $m\ge1$.
\writeup Fix an outcome $\omega$ lying in exactly $m$ of the sets. The three quantities $\sum_i\mathbf 1_{A_i}-\sum_{i<j}\mathbf 1_{A_i\cap A_j}$, $\mathbf 1_{\bigcup A_i}$ and $\sum_i\mathbf 1_{A_i}$ take at $\omega$ the values $m-\binom m2$, $\mathbf 1_{m\ge1}$ and $m$. For $m=0$ all three are $0$. For $m\ge1$, $m\ge1$ and $m-\binom m2=\frac{m(3-m)}2\le1$ (it equals $1$ for $m=1,2$ and is $\le0$ for $m\ge3$). So the inequalities hold between functions on $\Omega$, and taking expectations (which preserve inequalities and are linear, with $\mathbb E[\mathbf 1_A]=P(A)$) gives the stated ones. For the application, $P(\bigcup A_i)\le\sum P(A_i)\le n\cdot\frac1{2n}=\frac12$, so $P(\text{none})=1-P(\bigcup A_i)\ge\frac12$. Independence was not used; the union bound is the crudest and most used inequality of the probabilistic method, precisely because it asks nothing of the events.
\end{solution}

\begin{solution}{ex:records}
\idea A record at $i$ depends only on which of the first $i$ positions holds their maximum, and by symmetry that position is uniform. For pairs, the key is that conditioning on ``the maximum of the first $j$ is at $j$'' leaves the relative order of the first $j-1$ values uniform, so a record at $i<j$ is unaffected. With pairwise independence the variance is a sum, and Chebyshev finishes.
\writeup (a) $R_i=1$ iff the largest of $a_1,\dots,a_i$ sits at position $i$. For $s\le i$, swapping the entries at positions $s$ and $i$ is a bijection of the set of permutations that exchanges the events ``largest of the first $i$ at $s$'' and ``\dots at $i$'', so these $i$ events are equally likely, and disjoint with union $\Omega$: $P(R_i=1)=\frac1i$. By linearity $\mathbb ER=\sum_{i=1}^n\frac1i=H_n$.

(b) Let $i<j$. Count the permutations with records at $i$ and $j$: choose the set of values occupying positions $1,\dots,j$ ($\binom nj$ ways); put its largest at position $j$; arrange the remaining $j-1$ values in positions $1,\dots,j-1$ so that the largest of the first $i$ is at $i$, which by (a) applied to $j-1$ positions can be done in $\frac{(j-1)!}i$ ways; arrange the other $n-j$ values arbitrarily. Total $\binom nj\frac{(j-1)!}i(n-j)!=\frac{n!}{ij}$, so $P(R_i=R_j=1)=\frac1{ij}=P(R_i=1)P(R_j=1)$. Two indicators with $P(R_i=R_j=1)=P(R_i=1)P(R_j=1)$ are independent (the other three joint probabilities follow by complements), so Proposition~\ref{prop:varsum}(c) applies: $\operatorname{Var}R=\sum_i\operatorname{Var}R_i=\sum_i\frac1i\big(1-\frac1i\big)=H_n-\sum_i\frac1{i^2}$.

(c) Let $\mu=H_n$ and $v=\operatorname{Var}R$ with $n=10^6$. Write $\varepsilon_k=H_k-\ln k$. For $z=1/k>0$,
\[
\ln(1+z)-\frac z{1+z}<\frac{z^2}2,
\]
because $\varphi(z)=\frac{z^2}2-\ln(1+z)+\frac z{1+z}$ satisfies $\varphi(0)=0$ and $\varphi'(z)=z^2(2+z)/(1+z)^2\ge0$. Thus $\varepsilon_k-\varepsilon_{k+1}<1/(2k^2)$. Also $\ln(1+z)>z/(1+z)$, so $(\varepsilon_k)$ is strictly decreasing and $\varepsilon_n<\varepsilon_{20}$. The comparison $\sum_{k\ge20}1/k^2<\int_{19}^\infty x^{-2}\,\dd x=1/19$ gives
\[
\sum_{k\ge20}(\varepsilon_k-\varepsilon_{k+1})<\frac1{38},
\]
hence $\varepsilon_n>\varepsilon_{20}-1/38$ and
\[
\ln n+\varepsilon_{20}-\frac1{38}<\mu<\ln n+\varepsilon_{20}.
\]

The workshop of Chapter~\ref{ch:sequences} gives $0.693147<\ln2<0.693148$. The same geometric expansion on $[0,\frac14]$, stopped after eight terms, gives
\[
\ln\frac54=\sum_{k=1}^{8}\frac{(-1)^{k+1}}{k\,4^k}+r=\frac{12284087}{55050240}+r,
\]
with $0<r<1/(9\cdot4^9)$. These bounds are $0.223143<\ln(5/4)<0.223144$. Since $\ln10=3\ln2+\ln(5/4)$ and $\ln20=4\ln2+\ln(5/4)$,
\[
2.302584<\ln10<2.302588,\qquad 2.995731<\ln20<2.995736.
\]
The sum $H_{20}=\frac{55835135}{15519504}$ satisfies $3.59773<H_{20}<3.59774$, so $0.601994<\varepsilon_{20}<0.602009$. With $\ln n=6\ln10$,
\[
14.391<\mu<14.418.
\]

For the variance, $\sum_{i=1}^{20}1/i^2>1.5961$ and $\sum_{i=21}^n 1/i^2>\int_{21}^{n+1}x^{-2}\,\dd x>1/21-10^{-6}$. Therefore
\[
v<\mu-\bigl(1.5961+\tfrac1{21}-10^{-6}\bigr)<12.775.
\]
Chebyshev with $a=11.35$ gives $P(|R-\mu|\ge11.35)\le v/a^2<12.775/128.8225<0.1$. The interval $]\mu-11.35,\mu+11.35[$ lies inside $]3.04,25.77[\subset\,]3,26[$, so with probability greater than $0.9$ the integer $R$ satisfies $4\le R\le25$.
A million independent, identically distributed continuous observations have the same record count in law as this uniform permutation, because ties have probability zero and the ranks are uniform. Dependence or a changing distribution can alter the calculation.
\end{solution}

\begin{solution}{ex:poll}
\idea Each answer is a Bernoulli variable of parameter $p$; the poll is $n$ independent trials, and Theorem~\ref{thm:wlln} is exactly the statement asked. The population size is absent because sampling with replacement from a population is the same as repeating an experiment. For sharpness, the proof of Chebyshev loses only where $(X-\mathbb EX)^2$ exceeds $a^2$ by more than needed, so a variable that takes only the values $\mathbb EX\pm a$ and $\mathbb EX$ loses nothing.
\writeup (a) By Bernoulli's theorem with $\varepsilon=0.03$, $P(|\overline X_n-p|\ge0.03)\le\frac1{4n\cdot0.0009}=\frac{2500}{9n}\le\frac{277.8}n$, which is $\le0.05$ as soon as $n\ge5556$. (A normal approximation, proved in Volume~4, replaces the failure probability at level $0.05$ by a threshold near $1.96$ standard deviations. For $p=\frac12$ and margin $0.03$ that asks for about $1070$ people, the figure one reads in newspapers. Chebyshev, proved here from the variance alone, asks for $5556$.)

(b) The answers are independent Bernoulli($p$) variables whatever the population size, because sampling with replacement is independent repetition; only $p$ and $n$ enter. Without replacement the count is hypergeometric, and Proposition~\ref{prop:laws}(c) gives the variance $\frac{p(1-p)}n\cdot\frac{N-n}{N-1}$ for $\overline X_n$: smaller, so the bound only improves, and for $N$ in the millions and $n$ in the thousands the factor is indistinguishable from $1$. A poll of a thousand people is as good for China as for Luxembourg.

(c) Let $X$ take the values $-a$, $0$, $a$ with probabilities $\frac t2$, $1-t$, $\frac t2$, for some $t\in(0,1]$. Then $\mathbb EX=0$, $\operatorname{Var}X=a^2t$, and $P(|X|\ge a)=t=\frac{\operatorname{Var}X}{a^2}$: equality in Chebyshev. (Chebyshev's inequality is therefore the best possible statement that uses only the mean and the variance.)
\end{solution}

\begin{solution}{ex:ballot}
\idea Translate the count into a lattice path and the question into ``the path stays strictly above the axis''. The paths that fail are of two kinds: those that start with a vote for $B$ (easy to count) and those that start with a vote for $A$ but later touch the axis; the reflection argument pairs the second kind with the first, and the rest is arithmetic with binomial coefficients.
\writeup If $b=0$, every counting order is good and the formula gives $1$. Assume now $b\ge1$. A counting order is a sequence of $a$ steps $+1$ and $b$ steps $-1$; all $\binom{a+b}a$ sequences are equally likely (Proposition~\ref{prop:counts}(b), or counting paths as in Proposition~\ref{prop:paths}), and the partial sums form a path from $(0,0)$ to $(a+b,a-b)$. The count is \emph{good} if all partial sums are $>0$. A bad path either starts with $-1$ (there are $\binom{a+b-1}a$ of these: place the $a$ up-steps among the remaining $a+b-1$ positions) or starts with $+1$ and touches $0$ later. For a path of the second kind, let $t$ be the first time it hits $0$, and reflect its first $t$ steps (change their signs): the result starts with $-1$, hits $0$ at time $t$ for the first time, and continues as before. Conversely every path starting with $-1$ hits $0$ (it ends at $a-b>0$), and reflecting it up to its first return gives a path starting with $+1$ that touches $0$. The two operations are inverse, so the bad paths of the second kind are as many as those of the first, and
\[
P(\text{good})=1-\frac{2\binom{a+b-1}a}{\binom{a+b}a}=1-\frac{2b}{a+b}=\frac{a-b}{a+b},
\]
using $\binom{a+b-1}a=\frac{(a+b-1)!}{a!\,(b-1)!}=\frac b{a+b}\binom{a+b}a$. For $a=3$, $b=2$: one count in five, which the reader can verify by listing the $10$ orders.
\end{solution}

\begin{solution}{ex:ramsey}
\idea A random colouring has an expected number of monochromatic $k$-sets which is easy to compute by the indicator trick; if that expectation is below $1$, some colouring has none. The argument never says which, and this is the probabilistic method in its purest form: existence by averaging.
\writeup (a) For each $k$-set $S$ of vertices, let $I_S$ be the indicator that all $\binom k2$ edges inside $S$ have the same colour. All red has probability $2^{-\binom k2}$, all blue the same, so $\mathbb E[I_S]=2^{1-\binom k2}$, and by linearity the expected number of monochromatic $k$-sets is $\binom nk2^{1-\binom k2}$.

(b) With $n=\lfloor2^{k/2}\rfloor$, $\binom nk\le\frac{n^k}{k!}\le\frac{2^{k^2/2}}{k!}$, so the expectation is at most $\frac{2^{k^2/2+1-k(k-1)/2}}{k!}=\frac{2^{1+k/2}}{k!}$, which is $\frac{2^{5/2}}6<1$ for $k=3$ and decreases with $k$ (the ratio of consecutive terms is $\frac{\sqrt2}{k+1}<1$). A non-negative integer-valued random variable with expectation $<1$ takes the value $0$ with positive probability (otherwise it would be $\ge1$ everywhere and its expectation $\ge1$): some colouring has no monochromatic $k$-set. By definition $R(k,k)$ is the least $N$ such that every two-colouring of the complete graph on $N$ vertices has a monochromatic $k$-set, so $R(k,k)>n=\lfloor2^{k/2}\rfloor$, hence $R(k,k)>2^{k/2}$.

(c) The expectation is an average over all $2^{\binom n2}$ colourings, and an average below $1$ guarantees a term equal to $0$ without locating it; finding one explicitly is a different and much harder problem (challenge (2)). For $k=3$: on five vertices placed on a circle, colour the five sides of the pentagon red and the five diagonals blue; a triangle uses three of the ten edges, and no three sides form a triangle, no three diagonals either (the diagonals form a pentagram, which has no triangle), so $R(3,3)>5$. On six vertices, pick a vertex $v$: of its five edges, three have the same colour, say red, going to $x,y,z$. If one of the edges $xy,yz,xz$ is red it closes a red triangle with $v$; otherwise $xyz$ is a blue triangle. So $R(3,3)=6$.
\end{solution}

\begin{solution}{ex:covariance-sampling}
\idea Fixed sample size is an invariant: the indicators must sum to $k$.
That predicts a direction of zero variance before any matrix calculation.
Symmetry makes all diagonal covariance entries equal and all off-diagonal
entries equal, so only two scalar counts are needed. The resulting matrix is
a combination of $I$ and $J$, whose actions are transparent.
\writeup If $k=0$ or $k=n$, all indicators are deterministic, so $\Sigma=0$,
all eigenvalues are zero and the rank is zero. Otherwise counting subsets gives
\[
 P(X_i=1)=\frac{\binom{n-1}{k-1}}{\binom nk}=\frac kn=p.
\]
For $i\ne j$, counting gives
$P(X_i=X_j=1)=k(k-1)/(n(n-1))$; it is zero for $k=1$, as the same formula
states. Consequently
\[
 \operatorname{Var}X_i=p(1-p),\qquad
 \operatorname{Cov}(X_i,X_j)
 =\frac{k(k-1)}{n(n-1)}-\frac{k^2}{n^2}
 =-\frac{p(1-p)}{n-1}.
\]
These are exactly the entries of $p(1-p)(nI-J)/(n-1)$.
Writing $\boldsymbol1=(1,\ldots,1)^{\mathsf T}$, we have
$J\boldsymbol1=n\boldsymbol1$, and $Ju=0$ whenever $\sum_i u_i=0$.
Every vector has a unique decomposition into these two subspaces, by subtracting
its coordinate mean. The eigenvalues of $\Sigma$ are therefore $0$ on
$\R\boldsymbol1$ and $np(1-p)/(n-1)$ on the $(n-1)$-dimensional
sum-zero subspace. In the nondegenerate case $0<k<n$, its rank is $n-1$.
The zero variance says exactly that $\sum_iX_i=k$. Negative covariance says
that choosing one element leaves fewer places for the others: explicitly
$P(X_j=1\mid X_i=1)=(k-1)/(n-1)<k/n$.
\end{solution}

\begin{solution}{ex:arithmetic-independence}
\idea Divisibility by several distinct primes becomes one divisibility condition,
which can be counted exactly. Inclusion--exclusion compares the dependent finite
model with an independent product one intersection at a time. Keep the number
of error terms visible instead of dropping the floors informally.
\writeup For $I\subseteq\{1,\ldots,r\}$ put $d_I=\prod_{i\in I}p_i$,
with $d_\varnothing=1$. Unique prime factorisation from
Chapter~\ref{ch:sums} gives
\[
 P(p_i\mid U\text{ for every }i\in I)
 =\frac{\lfloor N/d_I\rfloor}{N}.
\]
If $q\mid N$, every $d_I\mid N$, so this probability is
$1/d_I=\prod_{i\in I}(1/p_i)$, proving mutual independence, including all
subfamilies. Inclusion--exclusion for the complement gives
\[
 P(\gcd(U,q)=1)
 =\sum_{I\subseteq\{1,\ldots,r\}}(-1)^{|I|}
       \frac{\lfloor N/d_I\rfloor}{N}.
\]
Replacing each floor by $N/d_I$ gives precisely
$\prod_i(1-1/p_i)$ by distributing the finite product. The empty-set term
has no error; each of the other $2^r-1$ terms has absolute error less than
$1/N$. The triangle inequality proves the stated bound (indeed it is strict).
For $p_1=2,p_2=3,N=3$, the two divisibility events each have probability $1/3$
and their intersection is empty, so they are not independent. For fixed primes,
the error bound tends to zero: the limiting density is an independent product
even though a particular finite sample need not have independent divisibility
events. This assertion keeps the set of primes fixed; the bound gives no uniform
control when $r$ grows without restriction.
\end{solution}

\subsection*{Notes on the challenges}

(1) The m\'enages problem is in \'Edouard Lucas's \emph{Th\'eorie des nombres} of 1891, and Touchard's formula was proved in his note of 1934; the inclusion--exclusion runs over the forbidden adjacencies, and the difficulty is to count the ways to choose $k$ pairwise compatible ``spouse next to spouse'' events around a circle, which is where the coefficient $\frac{2n}{2n-k}\binom{2n-k}k$ comes from (for the fixed-women convention of the challenge, it counts the $k$-subsets of a $2n$-cycle with no two adjacent, by a stars-and-bars argument around a circle). Bogart and Doyle gave a short non-sexist proof in 1986.

(2) $R(4,4)=18$ is from Greenwood and Gleason (1955), who used a construction with quadratic residues modulo $17$ for the lower bound. For $R(5,5)$ the bounds are $43$ (Exoo, 1989) and $46$ (Angeltveit and McKay, preprint 2024, published 2026). Erd\H os's 1947 bound $R(k,k)>2^{k/2}$ has been improved only in lower-order factors; matching an exponential lower bound $2^{ck}$, for a fixed $c>0$, by a deterministic explicit construction, with edge colours computable in time polynomial in $k$ from the endpoint labels, remains a central open problem. Erd\H os's story about the aliens is told by Joel Spencer in his \emph{Ten Lectures on the Probabilistic Method}.

(3) The finite tail-sum identity gives $\mathbb E[T]=\sum_{k=0}^NP(T>k)$, since $T\le N+1$. For $0\le k\le N$, $P(T>k)=\prod_{j<k}(1-\frac jN)\approx e^{-k^2/2N}$; the sum $\sum_ke^{-k^2/2N}\approx\int_0^\infty e^{-x^2/2N}\,dx=\sqrt{\pi N/2}$ by Exercise~\ref{ex:gaussian}. The integral suggests the leading scale but does not prove the refinement. Obtaining the constant $\frac23$ requires further terms together with uniform bounds on the tails; the refined asymptotic is a referenced outlook (see Flajolet and Sedgewick), not a complete exercise solution here. Ramanujan found the same constant in his $Q$-function in 1911, before anybody had asked about birthdays.

\begin{solution}{pb:cauchydav}
\idea If the sumset were small, a polynomial of controlled degree would vanish on $A\times B$; its key coefficient is a binomial
coefficient below $p$, hence non-zero.

\writeup (a) If $|A|+|B|>p$, then for every $c$ the sets $A$ and $c-B$ have total size $>p$, so they meet: $c\in A+B$. Suppose
$|A|+|B|\le p$ and $|A+B|\le|A|+|B|-2$, and choose $C\supset A+B$ with $|C|=|A|+|B|-2$. The polynomial
$f(x,y)=\prod_{c\in C}(x+y-c)$ vanishes on $A\times B$ and has degree $|A|+|B|-2$. Its top part is $(x+y)^{|A|+|B|-2}$, in which
the coefficient of $x^{|A|-1}y^{|B|-1}$ is $\binom{|A|+|B|-2}{|A|-1}$, not divisible by $p$ since $|A|+|B|-2<p$. Problem~\ref{pb:nullstellensatz}(b)
with $t_1=|A|-1<|A|$ and $t_2=|B|-1<|B|$ gives $(a,b)\in A\times B$ with $f(a,b)\ne0$, that is $a+b\notin C$: a contradiction.

(b) Induction on $k$: $|A_1+\dots+A_k|\ge\min\bigl(p,|A_1+\dots+A_{k-1}|+|A_k|-1\bigr)$, and if the first set already has $p$
elements so does the sum.

(c) Let $A_i=\{0,a_i\}$, of size $2$. By (b), the set of subset sums $A_1+\dots+A_{p-1}$ has at least
$\min(p,2(p-1)-(p-1)+1)=p$ elements. With $p-2$ copies of $1$, the residue $p-1$ is not a subset sum, so $p-1$ terms are needed.

\tools Chapter~\ref{ch:sums} (a prime $p$ divides $\binom pk$ for $0<k<p$); Problem~\ref{pb:nullstellensatz}.
\leads Cauchy proved (a) in 1813 to fill a gap in Lagrange's four squares theorem, and Davenport rediscovered it in 1935. Vosper
(1956) showed that equality with $|A|,|B|\ge2$ and $|A+B|\le p-2$ forces arithmetic progressions with the same difference. Over
$\Z$ and other groups the subject became additive combinatorics.
\end{solution}

\begin{solution}{pb:distances}
\idea Two circles meet in at most two points, so knowing two distances to two fixed points almost determines a point. Turned
around, two points have at most two common neighbours in the unit-distance graph, and Cauchy--Schwarz bounds graphs without
$K_{2,3}$.

\writeup (a) Let $t$ be the number of distinct distances, and fix $p\ne q$ in $P$. Every other point $x$ lies on one of the $t$ circles
centred at $p$ with a radius in the distance set, and on one of the $t$ such circles centred at $q$. Two circles with distinct
centres meet in at most $2$ points, so $n-2\le2t^2$.

(b) Let $G$ be the graph on $P$ joining points at distance $1$, with degrees $d_x$ and $u$ edges. Two distinct points $y,z$ have at
most $2$ common neighbours (the intersection of two unit circles). Counting the triples $(x,\{y,z\})$ with $y,z$ neighbours of $x$,
\[
\sum_x\binom{d_x}2\le2\binom n2.
\]
By convexity (Jensen, Theorem~\ref{thm:jensen}), $\sum_x\binom{d_x}2\ge n\binom{2u/n}2$, so $\frac{2u}n\bigl(\frac{2u}n-1\bigr)\le2(n-1)$. Solving the
quadratic in $\frac{2u}n$ gives $\frac{2u}n\le\frac12\bigl(1+\sqrt{8n-7}\bigr)$, which is the claim.

\tools Theorem~\ref{thm:jensen} and Theorem~\ref{thm:cs}; the K\H ov\'ari--S\'os--Tur\'an bound in Volume~4, for graphs with no $K_{s,t}$.
\leads Erdős's own bound was $\sqrt{n-\frac34}-\frac12$. The exponent $\frac12$ in (a) was raised many times, to $\frac45$ by Székely and
$0.864$ by Katz and Tardos, before Guth and Katz's $\frac{cn}{\log n}$ (2010). For unit distances, Problem~\ref{pb:szemereditrotter}(c)
improves (b) to $O(n^{4/3})$, the best bound known; the grid gives $n^{1+c/\log\log n}$, and Erdős conjectured this is the truth.
\end{solution}

\begin{solution}{pb:incidences}
\idea Lines through a point fan out: two points lie on at most one common line. So the number of point-pairs on lines is at most
$\binom n2$, which bounds $\sum k_\ell^2$, and Cauchy--Schwarz bounds $\sum k_\ell$.

\writeup (a) Let $k_\ell=|P\cap\ell|$, so $I=\sum_\ell k_\ell$. Each pair of points lies on at most one line of $L$, so
$\sum_\ell k_\ell(k_\ell-1)\le n(n-1)$, that is $\sum_\ell k_\ell^2\le n(n-1)+I$. By Cauchy--Schwarz, $I^2\le m\sum_\ell k_\ell^2\le m(n^2+I)$. So
$I^2-mI-mn^2\le0$, and $I\le\frac{m+\sqrt{m^2+4mn^2}}2\le\frac{m+m+2n\sqrt m}2=m+n\sqrt m$. Two lines meet in at most one point, so the same
argument with the roles of points and lines exchanged gives $I\le n+m\sqrt n$.

(b) Let $L$ be the set of the $\ell$ lines with at least $k$ points. Then $k\ell\le I(P,L)\le\ell\sqrt n+n$, so $\ell(k-\sqrt n)\le n$, and since
$k-\sqrt n\ge\frac k2$, $\ell\le\frac{2n}k$.

\tools Theorem~\ref{thm:cs}.
\leads The bound $n\sqrt m+m$ is sharp for the finite projective planes of Volume~9, where $n=m=q^2+q+1$ and every
line has $q+1$ points. In the real plane it is not: the Szemerédi--Trotter theorem (next problem) replaces $n\sqrt m$ by $(nm)^{2/3}$,
and the difference is a topological fact, the crossing lemma, that fails over finite fields.
\end{solution}

\begin{solution}{pb:szemereditrotter}
\idea Draw a graph whose vertices are the points and whose edges are the segments between consecutive points on each line. It
has many edges if there are many incidences, but few crossings, because two lines cross once. The crossing lemma says that
these two facts are incompatible unless the incidences are few.

\writeup (a) Discard the lines containing no point of $P$; this changes neither $I$ nor the bound. On each line $\ell$ with $k_\ell\ge1$
points, join consecutive points: $k_\ell-1$ edges. Two points lie on one line at most, so the graph $G$ on $P$ is simple, with
$e=\sum_\ell(k_\ell-1)\ge I-m$ edges. It is drawn with straight segments, and two edges cross only where their lines cross, so it has
at most $\binom m2<m^2$ crossings. If $e\ge4n$, the crossing lemma gives $\frac{e^3}{64n^2}\le m^2$, so $e\le4(nm)^{2/3}$ and
$I\le4(nm)^{2/3}+m$. If $e<4n$, then $I<4n+m$. In both cases $I\le4(nm)^{2/3}+m+4n$.

(b) Let $L$ be the set of the $\ell$ lines with at least $k$ points. By (a), $k\ell\le4(n\ell)^{2/3}+\ell+4n$. If $k\ge2$, then
$k\ell-\ell\ge\frac{k\ell}2$, so $\frac{k\ell}2\le4(n\ell)^{2/3}+4n$, and one of the two terms is at least $\frac{k\ell}4$. If $\frac{k\ell}4\le4(n\ell)^{2/3}$
then $\ell\le\frac{16^3n^2}{k^3}$; otherwise $\ell\le\frac{16n}k$.

(c) Let $C$ be the set of the $n$ unit circles centred at the points of $P$. The number $u$ of unit distances satisfies
$2u=I(P,C)$, the number of pairs (point, circle through it). Circles with at most $2$ points of $P$ contribute at most $2n$ incidences.
On each circle with $k\ge3$ points, join consecutive points along the circle by arcs: $k$ arcs. Two points are joined by at most
two arcs (they lie on at most two unit circles), so after removing duplicate arcs we get a simple graph with
$e\ge\frac{I-2n}2$ edges, drawn so that two arcs cross only where their circles cross: at most $2\binom n2<n^2$ crossings. If $e\ge4n$, the
crossing lemma gives $e^3\le64n^4$, $e\le4n^{4/3}$; otherwise $e<4n$. So $I\le2e+2n\le8n^{4/3}+10n$, and $u=\frac I2\le4n^{4/3}+5n$.

\tools the crossing lemma, proved in Volume~4; Problem~\ref{pb:incidences}.
\leads Szemerédi and Trotter proved (a) in 1983 with a much longer argument; Székely's proof appeared in 1997. The theorem is sharp
up to the constant: take the points of an $N\times2N^2$ grid and the lines $y=ax+b$ with $a\le N$, $b\le N^2$. Over $\mathbb F_q$ the
crossing lemma fails, and incidence bounds there are a current research topic.
\end{solution}

\begin{solution}{pb:sumproduct}
\idea If $A$ has few sums and few products, the grid $(A+A)\times(A\cdot A)$ is small, yet it carries $N^2$ lines with $N$ points each.
That is too many incidences for Szemerédi--Trotter.

\writeup Let $P=(A+A)\times(A\cdot A)$, with $|P|=|A+A|\,|A\cdot A|$, and for $b,c\in A$ let $\ell_{b,c}$ be the line $y=b(x-c)$. These $N^2$
lines are distinct ($b$ is the slope, $c$ the root). For $a\in A$, the point $(a+c,ab)$ lies on $\ell_{b,c}$ and in $P$, and these $N$ points
are distinct. So $I(P,L)\ge N^3$, and Problem~\ref{pb:szemereditrotter}(a) gives
\[
N^3\le4\bigl(|P|N^2\bigr)^{2/3}+N^2+4|P|.
\]
If $4|P|\ge\frac{N^3}3$, then $|P|\ge\frac{N^3}{12}\ge\frac{N^{5/2}}{12}$. Otherwise, for $N\ge3$, $N^2\le\frac{N^3}3$, so
$4(|P|N^2)^{2/3}\ge\frac{N^3}3$, that is $|P|\ge\frac{N^{5/2}}{12^{3/2}}$. For $N\le2$ the inequality $|P|\ge1\ge cN^{5/2}$ holds with $c=2^{-5/2}$.
So $|P|\ge cN^{5/2}$ with $c=12^{-3/2}$, and the larger of the two factors is at least $\sqrt c\,N^{5/4}$.

\tools Problem~\ref{pb:szemereditrotter}.
\leads Erdős and Szemerédi conjectured (1983) that $\max(|A+A|,|A\cdot A|)\ge N^{2-\varepsilon}$ for integers. Elekes's $\frac54$ was improved by
Solymosi to $\frac43$ (up to logarithms), by a clever count of slopes, and later slightly beyond $\frac43$ by Konyagin and Shkredov
and others. An arithmetic progression has few sums and a geometric progression few products: no set has both.
\end{solution}

\begin{solution}{pb:treebound}
\idea The matrix-tree theorem writes $\tau(G)$ as a product of eigenvalues whose sum is known; AM--GM bounds a product with a given
sum.

\writeup Let $0=\mu_1<\mu_2\le\dots\le\mu_n$ be the Laplacian eigenvalues ($\mu_2>0$ as $G$ is connected). By Kirchhoff's matrix-tree theorem in Volume~12, $\tau(G)=\frac1n\mu_2\cdots\mu_n$, and $\sum_i\mu_i=\operatorname{tr}L=\sum_vd_v=2m$. By AM--GM (Theorem~\ref{thm:amgm}),
\[
\mu_2\cdots\mu_n\le\Bigl(\frac{\mu_2+\dots+\mu_n}{n-1}\Bigr)^{n-1}=\Bigl(\frac{2m}{n-1}\Bigr)^{n-1},
\]
with equality if and only if $\mu_2=\dots=\mu_n=\mu$. In that case $L=\mu(I-\frac1nJ)$ (both sides vanish on $\mathbf1$ and are $\mu$ times the
identity on $\mathbf1^\perp$), so every off-diagonal entry of $L$ is $-\frac\mu n\ne0$: all edges are present, $G=K_n$. For $K_n$,
$m=\binom n2$ and the bound is $\frac1nn^{n-1}=n^{n-2}$: Cayley's formula.

\tools Kirchhoff's matrix-tree theorem in Volume~12; Theorem~\ref{thm:amgm}.
\leads For $d$-regular graphs this gives $\tau\le\frac1n\bigl(\frac{dn}{n-1}\bigr)^{n-1}$, roughly $d^n$; McKay (1983) found the sharp exponential rate
for large $d$-regular graphs, $\tau^{1/n}\to\frac{(d-1)^{d-1}}{(d^2-2d)^{d/2-1}}$, attained by random regular graphs.
\end{solution}

\begin{solution}{pb:parking}
\idea On a circle nobody fails to park and symmetry is perfect; cutting the circle at the empty spot gives the street. On the street, the first empty spot splits the
cars into a parking function and an unconstrained rest.

\writeup (a) On the circle of $n+1$ spots every car parks, since it drives around until it finds a free spot, and exactly one spot stays empty. Adding $c$ (mod $n+1$) to every
preference moves every car's spot by $c$, so it moves the empty spot by $c$. The $(n+1)^n$ preference sequences therefore fall into classes of $n+1$ rotations,
in each of which exactly one sequence leaves spot $n+1$ empty. A sequence leaves $n+1$ empty if and only if no car ever reaches spot $n+1$, that is, all its preferences are in
$\{1,\dots,n\}$ and every car parks in $1,\dots,n$ on the street: it is a parking function. So there are $\frac{(n+1)^n}{n+1}=(n+1)^{n-1}$.

(b) Suppose spot $k+1$ is the first empty one at the end. A car preferring a spot $\le k+1$ parks at the first free spot from its preference, which is at most $k+1$ since
$k+1$ stays free; it is not $k+1$, so it is in $\{1,\dots,k\}$. A car preferring a spot $\ge k+2$ parks beyond $k+1$. So the cars parked in $1,\dots,k$ are exactly the
$k$ cars with preference $\le k+1$; none prefers $k+1$ (it would park there), and their preferences form a parking function of length $k$. The other $n-k$ cars have
preferences in $\{k+2,\dots,n+1\}$, freely. Conversely such choices produce first empty spot $k+1$. This gives $\binom nk(k+1)^{k-1}(n-k)^{n-k}$ sequences, and summing over $k$ gives
$(n+1)^n$.

(c) $\binom nk\frac{(n-k)^{n-k}}{(n+1)^n}=\frac{n^{\underline k}}{n^kk!}\cdot\frac{(1-k/n)^{n-k}}{(1+1/n)^n}\to\frac{e^{-k}}{k!e}$. With $m=k+1$, $p_k=\frac{m^{m-2}e^{-m}}{(m-1)!}=\frac{m^{m-1}}{m!}e^{-m}$, so $\sum_kp_k=\sum_{m\ge1}\frac{m^{m-1}}{m!}e^{-m}$. The series
$T(x)$ has radius $\frac1e$ (Stirling), is continuous and increasing on $\left[0,\frac1e\right[$ with $T(0)=0$, and $T(x)e^{-T(x)}=x$. The map $t\mapsto te^{-t}$ increases from $0$ to $\frac1e$ on $[0,1]$, and
$T(x)$ never equals $1$ for $x<\frac1e$, so $T(x)<1$ and $T(x)$ is the inverse of $te^{-t}$ there; as $x\to\frac1e$, $T(x)\to1$. By monotone convergence, $\sum_m\frac{m^{m-1}}{m!}e^{-m}=1$.

(d) By Stirling, $\frac{m^{m-1}}{m!}e^{-m}\sim\frac{m^{m-1}e^{-m}}{\sqrt{2\pi m}\,m^me^{-m}}=\frac1{\sqrt{2\pi}m^{3/2}}$, and $m\sim k$. So $\sum_kkp_k=\infty$.

\tools{the tree-function identity $T=ze^T$; Theorem~\ref{thm:stirlingC}.}
\leads Konheim and Weiss (1966) introduced parking functions to analyse linear probing; the circular argument is Pollak's, recorded by Riordan (1969). The limit law
$(p_k)$ is the Borel distribution: $K+1$ has the law of the total size of a Galton--Watson tree with Poisson offspring of mean $1$, a critical tree, hence the heavy
tail. For finite $n$, one finds numerically $\mathbb E[K]\approx1.23\sqrt n$ at $n=3000$.
\end{solution}

\begin{solution}{pb:sinesecond}
\idea The quantity $\frac1{u_n^2}$ grows by almost exactly $\frac13$ at each step; the defect $\frac{u_n^2}{15}$ is about $\frac1{5n}$, so it adds up to $\frac15\ln n$. Inverting gives the second term
of $u_n$.

\writeup (a) $\sin u=u\bigl(1-\frac{u^2}6+\frac{u^4}{120}+O(u^6)\bigr)$, so $\sin^2u=u^2\bigl(1-\frac{u^2}3+\bigl(\frac1{36}+\frac1{60}\bigr)u^4+O(u^6)\bigr)=u^2\bigl(1-\varepsilon\bigr)$ with $\varepsilon=\frac{u^2}3-\frac{2u^4}{45}+O(u^6)$. Then
$\frac1{1-\varepsilon}=1+\varepsilon+\varepsilon^2+O(\varepsilon^3)=1+\frac{u^2}3+\bigl(\frac19-\frac2{45}\bigr)u^4+O(u^6)=1+\frac{u^2}3+\frac{u^4}{15}+O(u^6)$; divide by $u^2$.

(b) Let $e_k=\frac1{u_{k+1}^2}-\frac1{u_k^2}-\frac13$. The function $\frac1{\sin^2u}-\frac1{u^2}-\frac13$ is continuous on $\left]0,\frac\pi2\right]$ and $O(u^2)$ at $0$, so $|e_k|\le c\,u_k^2$. By Exercise~\ref{ex:sin}, $\frac1{u_{k+1}^2}-\frac1{u_k^2}\to\frac13$, so $\frac1{u_k^2}\ge\frac k4$ for $k\ge k_0$, and there is $c_1$ with $u_k^2\le\frac{c_1}{k+1}$ for all $k$. Hence $|\sum_{k<n}e_k|\le cc_1\sum_{k<n}\frac1{k+1}=O(\ln n)$ and
$\frac1{u_n^2}=\frac1{u_0^2}+\frac n3+\sum_{k<n}e_k=\frac n3+O(\ln n)$. Therefore $u_k^2=\frac3k+O(\frac{\ln k}{k^2})$, and by (a) $e_k=\frac{u_k^2}{15}+O(u_k^4)=\frac1{5k}+O(\frac{\ln k}{k^2})$. Summing,
$\sum_{k=1}^{n-1}e_k=\frac15\ln n+c'+o(1)$, and $\frac1{u_n^2}=\frac n3+\frac15\ln n+C+o(1)$.

(c) With $a=\frac15\ln n+C+o(1)$, $u_n=\bigl(\frac n3+a\bigr)^{-1/2}=\sqrt{\frac3n}\bigl(1+\frac{3a}n\bigr)^{-1/2}=\sqrt{\frac3n}\bigl(1-\frac{3a}{2n}+O(\frac{a^2}{n^2})\bigr)$, and $\sqrt{\frac3n}\cdot\frac3{2n}\cdot\frac{\ln n}5=\frac{3\sqrt3}{10}\cdot\frac{\ln n}{n^{3/2}}$.
The term in $C$ is $O(n^{-3/2})$.

\tools{Exercise~\ref{ex:sin} (the first term) and Chapter~\ref{ch:calculus} (Taylor expansions).}
\leads The correction is visible but slow: one finds numerically $\frac1{u_n^2}-\frac n3$ equal to $2.19$ at $n=10^3$ and $3.57$ at $n=10^6$, increasing by $\frac15\ln10=0.46$ for each factor $10$. For any
$u_{n+1}=u_n-au_n^p+bu_n^{2p-1}+\cdots$ the same computation gives a main term in $n^{-1/(p-1)}$ and a logarithmic correction, unless a coefficient happens to vanish.
\end{solution}

\section*{Sources and further reading}

\begin{itemize}\raggedright
\item R.~L.~Graham, D.~E.~Knuth and O.~Patashnik, \emph{Concrete Mathematics}, 2nd ed., Addison-Wesley, 1994, Chapters~5 and~8. Binomial coefficients, derangements, harmonic numbers and discrete probability, with a computer scientist's eye.
\item W.~Feller, \emph{An Introduction to Probability Theory and Its Applications}, Vol.~1, 3rd ed., Wiley, 1968, Chapters~II--IV and~IX. The matching problem (hats), the occupancy problems, inclusion--exclusion in probability, the ballot problem and the law of large numbers.
\item P.~Flajolet and R.~Sedgewick, \emph{Analytic Combinatorics}, Cambridge University Press, 2009 (free online). Cycles, fixed points and records of random permutations, through generating functions; the road from this chapter to Volume~4.
\item M.~Mitzenmacher and E.~Upfal, \emph{Probability and Computing}, 2nd ed., Cambridge University Press, 2017. The birthday problem, the coupon collector, the indicator trick and Chebyshev, applied to algorithms.
\item N.~Alon and J.~H.~Spencer, \emph{The Probabilistic Method}, 4th ed., Wiley, 2016; and J.~Spencer, \emph{Ten Lectures on the Probabilistic Method}, 2nd ed., SIAM, 1994. Exercise~\ref{ex:ramsey} is the first page of both.
\item F.~N.~David, \emph{Games, Gods and Gambling}, Griffin, 1962. The Pascal--Fermat letters of 1654 (translated in full), Huygens's \emph{De ratiociniis in ludo aleae} (1657), Montmort and the Bernoullis.
\item P.~R.~de Montmort, \emph{Essay d'analyse sur les jeux de hazard}, Paris, 1708; 2nd ed.\ 1713, with the correspondence with Nicolaus Bernoulli on the \emph{probl\`eme des rencontres}. L.~Euler, \emph{Calcul de la probabilit\'e dans le jeu de rencontre}, M\'emoires de l'Acad\'emie des sciences de Berlin 7 (1751, printed 1753), 255--270, for the series and the recurrence of Exercise~\ref{ex:derangrec}.
\item J.~Bernoulli, \emph{Ars Conjectandi}, Basel, 1713; English translation by E.~D.~Sylla, Johns Hopkins University Press, 2006. Part~IV contains the ``golden theorem'' of Theorem~\ref{thm:wlln}, with the twenty-year story in the translator's introduction.
\item P.~L.~Chebyshev, \emph{Des valeurs moyennes}, J.\ Math.\ Pures Appl.\ 12 (1867), 177--184; the inequality was published earlier by I.-J.~Bienaym\'e, C.\ R.\ Acad.\ Sci.\ Paris 37 (1853), 309--324, which is why it is often called the Bienaym\'e--Chebyshev inequality.
\item R.~von Mises, \emph{\"Uber Aufteilungs- und Besetzungs-\allowbreak Wahrscheinlichkeiten}, Revue de la Facult\'e des Sciences de l'Universit\'e d'Istanbul 4 (1939), 145--163. The birthday problem.
\item P.~Erd\H{o}s and G.~Szekeres, \emph{A combinatorial problem in geometry}, Compositio Math.\ 2 (1935), 463--470; P.~Erd\H os, \emph{Some remarks on the theory of graphs}, Bull.\ Amer.\ Math.\ Soc.\ 53 (1947), 292--294. Exercises~\ref{ex:erdos} and~\ref{ex:ramsey}. R.~E.~Greenwood and A.~M.~Gleason, Canad.\ J.\ Math.\ 7 (1955), 1--7, for $R(4,4)=18$; V.~Angeltveit and B.~D.~McKay, \emph{$R(5,5)\le46$}, Journal of Graph Theory (2026), doi:10.1002/jgt.70029; preprint arXiv:2409.15709 (2024).
\item S.~Selvin, \emph{A problem in probability}, Amer.\ Statist.\ 29 (1975), 67, where the Monty Hall problem was first posed; M.~vos Savant, \emph{Ask Marilyn}, Parade, 9 September 1990 and the following months, for the storm of letters.
\item K.~Barner, \emph{Neues zu Fermats Geburtsdatum}, Mitteilungen der Deutschen Mathematiker-Vereinigung 15 (2007), 12--14. Archival evidence for a birth in 1607 rather than the traditional 1601.
\item M.~Renault, \emph{Lost (and found) in translation: Andr\'e's actual method and its application to the generalized ballot problem}, Amer.\ Math.\ Monthly 115 (2008), 358--363.
\item J.~Bertrand, C.\ R.\ Acad.\ Sci.\ Paris 105 (1887), 369, and D.~Andr\'e, ibid., 436--437. The ballot problem and Andr\'e's original method.
\item C.~E.~Bonferroni, \emph{Teoria statistica delle classi e calcolo delle probabilit\`a}, Pubbl.\ R.\ Ist.\ Sup.\ Sci.\ Econ.\ Comm.\ Firenze 8 (1936), 3--62. A.~R\'enyi, \emph{Th\'eorie des \'el\'ements saillants d'une suite d'observations}, Ann.\ Fac.\ Sci.\ Univ.\ Clermont 8 (1962), 7--13, for records.
\item \'E.~Lucas, \emph{Th\'eorie des nombres}, Gauthier-Villars, 1891, p.~491; J.~Touchard, C.\ R.\ Acad.\ Sci.\ Paris 198 (1934), 631--633; K.~P.~Bogart and P.~G.~Doyle, \emph{Non-sexist solution of the m\'enage problem}, Amer.\ Math.\ Monthly 93 (1986), 514--519.
\item P.~G.~L.~Dirichlet, \emph{Verallgemeinerung eines Satzes aus der Lehre von den Kettenbr\"uchen}, Bericht Berlin Akad.\ (1842), 93--95. The pigeonhole principle applied to rational approximation.
\end{itemize}
